% 附录 F Third-Party Licenses
\biappendix{第三方许可}{Third-Party Licenses}
\label{app:f}

本附录提供 IC Validator 工具所用第三方软件的许可信息。
\begin{itemize}
  \item 许可概述（Licensing Overview）
  \item OSS 软件包声明（OSS Package Notices）
  \begin{itemize}
    \item ANTLR
    \item Boost
    \item cx\_Freeze
    \item Flatbuffers
    \item libdp
    \item libxml2
    \item MariaDB Connector
    \item MariaDB
    \item MCPP Public Domain Code
    \item NumPy*
    \item patchELF
    \item pyparsing*
    \item Python 3.6 License
    \item scikit-learn
    \item SciPy*
    \item Shroud-1.0
    \item SQLite
    \item TclLib
    \item Tcl/Tk（注：tbcload 代码是 Tcl 的一部分）
    \item zlib
    \item Zstandard
  \end{itemize}
  \item 标准 OSS 许可文本（Standard OSS License Text）
  \begin{itemize}
    \item Apache-2.0
    \item Artistic-1.0-Perl
    \item GPL-1.0
    \item GPL-3.0
  \end{itemize}
\end{itemize}

% ============================================================
\section{许可概述}
\subsection*{Licensing Overview}

本文档包含与 Synopsys\textsuperscript{\textregistered} IC Validator 产品（下称``SOFTWARE''）所含自由与开源软件（``OSS''）相关的许可信息。适用的 OSS 许可条款约束 Synopsys\textsuperscript{\textregistered} 对 SOFTWARE 的分发以及您对 SOFTWARE 的使用。Synopsys\textsuperscript{\textregistered} 以及 SOFTWARE 的第三方作者、许可方与分发方，对因任何使用与分发 SOFTWARE 而产生的全部保证与全部责任予以免责。若 OSS 在与 Synopsys\textsuperscript{\textregistered} 的协议下提供，且该协议不同于适用的 OSS 许可，则那些条款仅由 Synopsys\textsuperscript{\textregistered} 单独提供。

Synopsys\textsuperscript{\textregistered} 已在下文转载 OSS 软件包中出现的版权与其他许可声明。尽管 Synopsys\textsuperscript{\textregistered} 力求为每个 OSS 软件包提供完整准确的版权与许可信息，但 Synopsys\textsuperscript{\textregistered} 不声明或保证下列信息完整、正确或无误。鼓励 SOFTWARE 接收方：(a)~核查所识别的 OSS 软件包以确认本文所提供许可信息的准确性；(b)~将本文档中发现的任何不准确或错误通知 Synopsys\textsuperscript{\textregistered}，以便 Synopsys\textsuperscript{\textregistered} 相应更新本文档。

某些 OSS 许可（例如 GNU General Public Licenses、GNU Library/Lesser General Public Licenses、Affero General Public Licenses、Mozilla Public Licenses、Common Development and Distribution Licenses、Common Public License 以及 Eclipse Public License）要求：与所分发 OSS 二进制相对应的源代码，须按同一 OSS 许可的条款向接收方或其他请求者提供。

希望收到此类对应源代码副本的接收方或请求者，应向 Synopsys\textsuperscript{\textregistered} 邮寄请求至：
\begin{quote}
Synopsys\\
Attn: Open Source Requests\\
690 E. Middlefield Road\\
Mountain View, CA 94043
\end{quote}

提交的全部 OSS 请求中请提供下列信息：
\begin{itemize}
  \item 您请求源代码的 OSS 软件包；
  \item 所请求 OSS 软件包随其分发的 Synopsys\textsuperscript{\textregistered} 产品（以及任何可用的版本信息）；
  \item Synopsys\textsuperscript{\textregistered} 可就该请求联系您的电子邮箱（如有）；以及
  \item 请求源代码的投递邮寄地址。
\end{itemize}

已在 OSS 软件包名称后添加星号（*），以表示这些组件可能包含在使用 ZeBu 开发的软件或硬件产品中。这些组件的 OSS 许可信息也转载于随 SOFTWARE 提供的 IC Validator 可再分发 OSS 声明文本文件中。

\begin{noteBox}
下列各 OSS 软件包与标准许可小节中的法律许可原文保留英文；中文仅译章节说明与概述。
\end{noteBox}

% ============================================================
\section{OSS 软件包声明}
\subsection*{OSS Package Notices}

\subsection{ANTLR}
\subsubsection*{ANTLR}

以下为 ANTLR 的许可声明原文。

\begin{lstlisting}[basicstyle=\ttfamily\footnotesize,breaklines=true]
     ANTLR 2 License
     [The BSD License]
     Copyright (c) 2012 Terence Parr and Sam Harwell
     All rights reserved.




     Redistribution and use in source and binary forms, with or without modification, are
     permitted provided that the following conditions are met:
      • Redistributions of source code must retain the above copyright notice, this list of
        conditions and the following disclaimer.
      • Redistributions in binary form must reproduce the above copyright notice, this list of
        conditions and the following disclaimer in the documentation and/or other materials
        provided with the distribution.
      • Neither the name of the author nor the names of its contributors may be used to
        endorse or promote products derived from this software without specific prior written
        permission.
     THIS SOFTWARE IS PROVIDED BY THE COPYRIGHT HOLDERS AND
     CONTRIBUTORS "AS IS" AND ANY EXPRESS OR IMPLIED WARRANTIES,
     INCLUDING, BUT NOT LIMITED TO, THE IMPLIED WARRANTIES OF
     MERCHANTABILITY AND FITNESS FOR A PARTICULAR PURPOSE ARE
     DISCLAIMED. IN NO EVENT SHALL THE COPYRIGHT OWNER OR CONTRIBUTORS
     BE LIABLE FOR ANY DIRECT, INDIRECT, INCIDENTAL, SPECIAL, EXEMPLARY, OR
     CONSEQUENTIAL DAMAGES (INCLUDING, BUT NOT LIMITED TO, PROCUREMENT
     OF SUBSTITUTE GOODS OR SERVICES; LOSS OF USE, DATA, OR PROFITS;
     OR BUSINESS INTERRUPTION) HOWEVER CAUSED AND ON ANY THEORY OF
     LIABILITY, WHETHER IN CONTRACT, STRICT LIABILITY, OR TORT (INCLUDING
     NEGLIGENCE OR OTHERWISE) ARISING IN ANY WAY OUT OF THE USE OF THIS
     SOFTWARE, EVEN IF ADVISED OF THE POSSIBILITY OF SUCH DAMAGE.
\end{lstlisting}

\subsection{Boost}
\subsubsection*{Boost}

以下为 Boost 的许可声明原文。

\begin{lstlisting}[basicstyle=\ttfamily\footnotesize,breaklines=true]
     Boost Software License - Version 1.0 - August 17th, 2003
     Permission is hereby granted, free of charge, to any person or organization obtaining
     a copy of the software and accompanying documentation covered by this license (the
     “Software”) to use, reproduce, display, distribute, execute, and transmit the Software,
     and to prepare derivative works of the Software, and to permit third-parties to whom the
     Software is furnished to do so, all subject to the following:
     The copyright notices in the Software and this entire statement, including the above
     license grant, this restriction and the following disclaimer, must be included in all copies
     of the Software, in whole or in part, and all derivative works of the Software, unless such
     copies or derivative works are solely in the form of machine-executable object code
     generated by a source language processor.
     THE SOFTWARE IS PROVIDED “AS IS”, WITHOUT WARRANTY OF ANY KIND,
     EXPRESS OR IMPLIED, INCLUDING BUT NOT LIMITED TO THE WARRANTIES OF
     MERCHANTABILITY, FITNESS FOR A PARTICULAR PURPOSE, TITLE AND NON-
     INFRINGEMENT. IN NO EVENT SHALL THE COPYRIGHT HOLDERS OR ANYONE




     DISTRIBUTING THE SOFTWARE BE LIABLE FOR ANY DAMAGES OR OTHER
     LIABILITY, WHETHER IN CONTRACT, TORT OR OTHERWISE, ARISING FROM, OUT
     OF OR IN CONNECTION WITH THE SOFTWARE OR THE USE OR OTHER DEALINGS
     IN THE SOFTWARE.
\end{lstlisting}

\subsection{cx\_Freeze}
\subsubsection*{cx\_Freeze}

以下为 cx\_Freeze 的许可声明原文。

\begin{lstlisting}[basicstyle=\ttfamily\footnotesize,breaklines=true]
     Project homepage/download site: http://cx-freeze.readthedocs.io/; https://pypi.python.org/
     pypi/cx_Freeze/4.3.4
     Project licensing notices: http://cx-freeze.readthedocs.io/en/latest/license.html
     Copyright © 2007-2017, Anthony Tuininga.
     Copyright © 2001-2006, Computronix (Canada) Ltd., Edmonton, Alberta,
      Canada.
     All rights reserved.

     NOTE: this license is derived from the Python Software Foundation License
     which can be found at http://www.python.org/psf/license

     License for cx_Freeze 6.0b1

     1. This LICENSE AGREEMENT is between the copyright holders and the
     Individual or Organization ("Licensee") accessing and otherwise
     using cx_Freeze software in source or binary form and its associated
     documentation.

     2. Subject to the terms and conditions of this License Agreement,
     the copyright holders hereby grant Licensee a nonexclusive, royalty-
     free, world-wide license to reproduce, analyze, test, perform and/or
     display publicly, prepare derivative works, distribute, and otherwise
     use cx_Freeze alone or in any derivative version, provided, however,
     that this License Agreement and this notice of copyright are retained in
     cx_Freeze alone or in any derivative version prepared by Licensee.

     3. In the event Licensee prepares a derivative work that is based on or incorporates
     cx_Freeze or any part thereof, and wants to make the derivative work available to others
     as provided herein, then Licensee hereby agrees to include in any such work a brief
     summary of the changes made to cx_Freeze.
     4. The copyright holders are making cx_Freeze available to Licensee on an “AS IS” basis.
     THE COPYRIGHT HOLDERS MAKE NO REPRESENTATIONS OR WARRANTIES,
     EXPRESS OR IMPLIED. BY WAY OF EXAMPLE, BUT NOT LIMITATION, THE
     COPYRIGHT HOLDERS MAKE NO AND DISCLAIM ANY REPRESENTATION OR
     WARRANTY OF MERCHANTABILITY OR FITNESS FOR ANY PARTICULAR PURPOSE
     OR THAT THE USE OF CX_FREEZE WILL NOT INFRINGE ANY THIRD PARTY
     RIGHTS.




     5. THE COPYRIGHT HOLDERS SHALL NOT BE LIABLE TO LICENSEE OR
     ANY OTHER USERS OF CX_FREEZE FOR ANY INCIDENTAL, SPECIAL,
     OR CONSEQUENTIAL DAMAGES OR LOSS AS A RESULT OF MODIFYING,
     DISTRIBUTING, OR OTHERWISE USING CX_FREEZE, OR ANY DERIVATIVE
     THEREOF, EVEN IF ADVISED OF THE POSSIBILITY THEREOF.
     6. This License Agreement will automatically terminate upon a material breach of its terms
     and conditions.
     7. Nothing in this License Agreement shall be deemed to create any relationship of
     agency, partnership, or joint venture between the copyright holders and Licensee. This
     License Agreement does not grant permission to use copyright holder’s trademarks or
     trade name in a trademark sense to endorse or promote products or services of Licensee,
     or any third party.
     8. By copying, installing or otherwise using cx_Freeze, Licensee agrees to be bound
     by the terms and conditions of this License Agreement. Computronix® is a registered
     trademark of Computronix (Canada) Ltd.
\end{lstlisting}

\subsection{Flatbuffers}
\subsubsection*{Flatbuffers}

以下为 Flatbuffers 的许可声明原文。

\begin{lstlisting}[basicstyle=\ttfamily\footnotesize,breaklines=true]
     Project homepage/download site: http://google.github.io/flatbuffers/; https://github.com/
     google/flatbuffers
     Project licensing notices: https://github.com/google/flatbuffers/blob/master/LICENSE.txt
     LICENSE.txt:
     See Apache-2.0 in the Standard OSS License Text on page 542."
\end{lstlisting}

\subsection{libdp}
\subsubsection*{libdp}

以下为 libdp 的许可声明原文。

\begin{lstlisting}[basicstyle=\ttfamily\footnotesize,breaklines=true]
      1. The "Software", below, refers to the Tcl-DP system, developed by the Tcl-DP group (in
         either source-code, object-code or executable-code form), and related documentation,
         and a "work based on the Software" means a work based on either the Software, on
         part of the Software, or on any derivative work of the Software under copyright law:
         that is, a work containing all or a portion of the Tcl-DP system, either verbatim or with
         modifications. Each licensee is addressed as "you" or "Licensee."
      2. Cornell University as the parent organization of the Tcl-DP group holds copyrights in
         the Software. The copyright holder reserves all rights except those expressly granted to
         licensees, and U.S. Government license rights.
      3. Permission is hereby granted to use, copy, modify, and to redistribute to others. If
         you distribute a copy or copies of the Software, or you modify a copy or copies of the




         Software or any portion of it, thus forming a work based on the Software, and make
         and/or distribute copies of such work, you must meet the following conditions:
             1. If you make a copy of the Software (modified or verbatim) it must include the
             copyright notice and this license.
             2. You must cause the modified Software to carry prominent notices stating that you
             changed specified portions of the Software.
      4. LICENSEE AGREES THAT THE EXPORT OF GOODS AND/OR TECHNICAL DATA
         FROM THE UNITED STATES MAY REQUIRE SOME FORM OF EXPORT CONTROL
         LICENSE FROM THE U.S. GOVERNMENT AND THAT FAILURE TO OBTAIN SUCH
         EXPORT CONTROL LICENSE MAY RESULT IN CRIMINAL LIABILITY UNDER U.S.
         LAWS.
      5. Portions of the Software resulted from work developed under a U.S. Government
         Contract and are subject to the following license: the Government is granted for itself
         and others acting on its behalf a paid-up, nonexclusive, irrevocable worldwide license
         in this computer software to reproduce, prepare derivative works, and perform publicly
         and display publicly.
      6. Disclaimer of warranty: Licensor provides the software on an ``as is'' basis. Licensor
         does not warrant, guarantee, or make any representations regarding the use or results
         of the software with respect to its correctness, accuracy, reliability or performance.
         The entire risk of the use and performance of the software is assumed by licensee.
         ALL WARRANTIES INCLUDING, WITHOUT LIMITATION, ANY WARRANTY OF
         FITNESS FOR A PARTICULAR PURPOSE OR MERCHANTABILITY ARE HEREBY
         EXCLUDED.
      7. Lack of maintenance or support services: Licensee understands and agrees that
         licensor is under no obligation to provide maintenance, support or update services,
         notices of latent defects, or correction of defects for the software.
      8. Limitation of liability, indemnification: Even if advised of the possibility of damages,
         under no circumstances shall licensor be liable to licensee or any third party for
         damages of any character, including, without limitation, direct, indirect, incidental,
         consequential or special damages, loss of profits, loss of use, loss of goodwill,
         computer failure or malfunction. Licensee agrees to indemnify and hold harmless
         licensor for any and all liability licensor may incur as a result of licensee's use of the
         software.
\end{lstlisting}

\subsection{libxml2}
\subsubsection*{libxml2}

以下为 libxml2 的许可声明原文。

\begin{lstlisting}[basicstyle=\ttfamily\footnotesize,breaklines=true]
     Libxml2 is the XML C parser and toolkit developed for the Gnome project (but usable
     outside of the Gnome platform), it is free software available under the MIT License. The
     MIT license:
     The MIT License (MIT)
     Copyright (c) <year> <copyright holders>
     Permission is hereby granted, free of charge, to any person obtaining a copy of this
     software and associated documentation files (the "Software"), to deal in the Software
     without restriction, including without limitation the rights to use, copy, modify, merge,
     publish, distribute, sublicense, and/or sell copies of the Software, and to permit persons to
     whom the Software is furnished to do so, subject to the following conditions:
     The above copyright notice and this permission notice shall be included in all copies or
     substantial portions of the Software.
     THE SOFTWARE IS PROVIDED "AS IS", WITHOUT WARRANTY OF ANY KIND,
     EXPRESS OR IMPLIED, INCLUDING BUT NOT LIMITED TO THE WARRANTIES
     OF MERCHANTABILITY, FITNESS FOR A PARTICULAR PURPOSE AND
     NONINFRINGEMENT. IN NO EVENT SHALL THE AUTHORS OR COPYRIGHT
     HOLDERS BE LIABLE FOR ANY CLAIM, DAMAGES OR OTHER LIABILITY, WHETHER
     IN AN ACTION OF CONTRACT, TORT OR OTHERWISE, ARISING FROM, OUT OF OR
     IN CONNECTION WITH THE SOFTWARE OR THE USE OR OTHER DEALINGS IN THE
     SOFTWARE.
\end{lstlisting}

\subsection{MariaDB Connector}
\subsubsection*{MariaDB Connector}

以下为 MariaDB Connector 的许可声明原文。

\begin{lstlisting}[basicstyle=\ttfamily\footnotesize,breaklines=true]
     Synopsys modifications: NoneMariaDB Connector is licensed under the GNU Lesser
     General Public License, version 2.1, (LGPL-2.1). Software licensed under the LGPL-2.1
     is free software that comes with ABSOLUTELY NO WARRANTY. You may modify,
     redistribute, or otherwise use the GPL software under the terms of the LGPL-2.1 (http://
     www.gnu.org/licenses/old-licenses/lgpl-2.1.html).
     GNU LESSER GENERAL PUBLIC LICENSE
     Version 2.1, February 1999

     Copyright (C) 1991, 1999 Free Software Foundation, Inc. 51 Franklin
     Street, Fifth Floor, Boston, MA 02110-1301 USA Everyone is permitted
     to copy and distribute verbatim copies of this license document, but
     changing it is not allowed.




     [This is the first released version of the Lesser GPL. It also counts as
     the successor of the GNU Library Public License, version 2, hence the
     version number 2.1.]

     TERMS AND CONDITIONS FOR COPYING, DISTRIBUTION AND MODIFICATION
     0. This License Agreement applies to any software library or other program which
     contains a notice placed by the copyright holder or other authorized party saying it may
     be distributed under the terms of this Lesser General Public License (also called "this
     License"). Each licensee is addressed as "you".
     A "library" means a collection of software functions and/or data prepared so as to be
     conveniently linked with application programs (which use some of those functions and
     data) to form executables.
     The "Library", below, refers to any such software library or work which has been distributed
     under these terms. A "work based on the Library" means either the Library or any
     derivative work under copyright law: that is to say, a work containing the Library or a
     portion of it, either verbatim or with modifications and/or translated straightforwardly
     into another language. (Hereinafter, translation is included without limitation in the term
     "modification".)
     "Source code" for a work means the preferred form of the work for making modifications
     to it. For a library, complete source code means all the source code for all modules it
     contains, plus any associated interface definition files, plus the scripts used to control
     compilation and installation of the library.
     Activities other than copying, distribution and modification are not covered by this License;
     they are outside its scope. The act of running a program using the Library is not restricted,
     and output from such a program is covered only if its contents constitute a work based on
     the Library (independent of the use of the Library in a tool for writing it). Whether that is
     true depends on what the Library does and what the program that uses the Library does.
      1. You may copy and distribute verbatim copies of the Library's complete source code
         as you receive it, in any medium, provided that you conspicuously and appropriately
         publish on each copy an appropriate copyright notice and disclaimer of warranty; keep
         intact all the notices that refer to this License and to the absence of any warranty; and
         distribute a copy of this License along with the Library.
         You may charge a fee for the physical act of transferring a copy, and you may at your
         option offer warranty protection in exchange for a fee.




      2. You may modify your copy or copies of the Library or any portion of it, thus forming a
         work based on the Library, and copy and distribute such modifications or work under
         the terms of Section 1 above, provided that you also meet all of these conditions:
          ◦ The modified work must itself be a software library.
          ◦ You must cause the files modified to carry prominent notices stating that you
            changed the files and the date of any change.
          ◦ You must cause the whole of the work to be licensed at no charge to all third parties
            under the terms of this License.
          ◦ If a facility in the modified Library refers to a function or a table of data to be
            supplied by an application program that uses the facility, other than as an argument
            passed when the facility is invoked, then you must make a good faith effort to
            ensure that, in the event an application does not supply such function or table, the
            facility still operates, and performs whatever part of its purpose remains meaningful.
         (For example, a function in a library to compute square roots has a purpose that is
         entirely well-defined independent of the application. Therefore, Subsection 2d requires
         that any application-supplied function or table used by this function must be optional:
         if the application does not supply it, the square root function must still compute square
         roots.)
         These requirements apply to the modified work as a whole. If identifiable sections
         of that work are not derived from the Library, and can be reasonably considered
         independent and separate works in themselves, then this License, and its terms, do
         not apply to those sections when you distribute them as separate works. But when
         you distribute the same sections as part of a whole which is a work based on the
         Library, the distribution of the whole must be on the terms of this License, whose
         permissions for other licensees extend to the entire whole, and thus to each and every
         part regardless of who wrote it.
         Thus, it is not the intent of this section to claim rights or contest your rights to work
         written entirely by you; rather, the intent is to exercise the right to control the distribution
         of derivative or collective works based on the Library.
         In addition, mere aggregation of another work not based on the Library with the Library
         (or with a work based on the Library) on a volume of a storage or distribution medium
         does not bring the other work under the scope of this License.
      3. You may opt to apply the terms of the ordinary GNU General Public License instead of
         this License to a given copy of the Library. To do this, you must alter all the notices that
         refer to this License, so that they refer to the ordinary GNU General Public License,
         version 2, instead of to this License. (If a newer version than version 2 of the ordinary
         GNU General Public License has appeared, then you can specify that version instead if
         you wish.) Do not make any other change in these notices.




         Once this change is made in a given copy, it is irreversible for that copy, so the ordinary
         GNU General Public License applies to all subsequent copies and derivative works
         made from that copy.
         This option is useful when you wish to copy part of the code of the Library into a
         program that is not a library.
      4. You may copy and distribute the Library (or a portion or derivative of it, under Section
         2) in object code or executable form under the terms of Sections 1 and 2 above
         provided that you accompany it with the complete corresponding machine-readable
         source code, which must be distributed under the terms of Sections 1 and 2 above on a
         medium customarily used for software interchange.
         If distribution of object code is made by offering access to copy from a designated
         place, then offering equivalent access to copy the source code from the same place
         satisfies the requirement to distribute the source code, even though third parties are
         not compelled to copy the source along with the object code.
      5. A program that contains no derivative of any portion of the Library, but is designed to
         work with the Library by being compiled or linked with it, is called a "work that uses the
         Library". Such a work, in isolation, is not a derivative work of the Library, and therefore
         falls outside the scope of this License.
         However, linking a "work that uses the Library" with the Library creates an executable
         that is a derivative of the Library (because it contains portions of the Library), rather
         than a "work that uses the library". The executable is therefore covered by this License.
         Section 6 states terms for distribution of such executables.
         When a "work that uses the Library" uses material from a header file that is part of
         the Library, the object code for the work may be a derivative work of the Library even
         though the source code is not. Whether this is true is especially significant if the work
         can be linked without the Library, or if the work is itself a library. The threshold for this
         to be true is not precisely defined by law.
         If such an object file uses only numerical parameters, data structure layouts and
         accessors, and small macros and small inline functions (ten lines or less in length),
         then the use of the object file is unrestricted, regardless of whether it is legally a
         derivative work. (Executables containing this object code plus portions of the Library
         will still fall under Section 6.)
         Otherwise, if the work is a derivative of the Library, you may distribute the object code
         for the work under the terms of Section 6. Any executables containing that work also
         fall under Section 6, whether or not they are linked directly with the Library itself.
      6. As an exception to the Sections above, you may also combine or link a "work that
         uses the Library" with the Library to produce a work containing portions of the Library,
         and distribute that work under terms of your choice, provided that the terms permit




         modification of the work for the customer's own use and reverse engineering for
         debugging such modifications.
         You must give prominent notice with each copy of the work that the Library is used in it
         and that the Library and its use are covered by this License. You must supply a copy of
         this License. If the work during execution displays copyright notices, you must include
         the copyright notice for the Library among them, as well as a reference directing the
         user to the copy of this License. Also, you must do one of these things:
          ◦ Accompany the work with the complete corresponding machine-readable source
            code for the Library including whatever changes were used in the work (which must
            be distributed under Sections 1 and 2 above); and, if the work is an executable
            linked with the Library, with the complete machine-readable "work that uses the
            Library", as object code and/or source code, so that the user can modify the Library
            and then relink to produce a modified executable containing the modified Library.
            (It is understood that the user who changes the contents of definitions files in the
            Library will not necessarily be able to recompile the application to use the modified
            definitions.)
          ◦ Use a suitable shared library mechanism for linking with the Library. A suitable
            mechanism is one that (1) uses at run time a copy of the library already present
            on the user's computer system, rather than copying library functions into the
            executable, and (2) will operate properly with a modified version of the library, if the
            user installs one, as long as the modified version is interface-compatible with the
            version that the work was made with.
          ◦ Accompany the work with a written offer, valid for at least three years, to give the
            same user the materials specified in Subsection 6a, above, for a charge no more
            than the cost of performing this distribution.
          ◦ If distribution of the work is made by offering access to copy from a designated
            place, offer equivalent access to copy the above specified materials from the same
            place.
          ◦ Verify that the user has already received a copy of these materials or that you have
            already sent this user a copy.
         For an executable, the required form of the "work that uses the Library" must include
         any data and utility programs needed for reproducing the executable from it. However,
         as a special exception, the materials to be distributed need not include anything that
         is normally distributed (in either source or binary form) with the major components
         (compiler, kernel, and so on) of the operating system on which the executable runs,
         unless that component itself accompanies the executable.
         It may happen that this requirement contradicts the license restrictions of other
         proprietary libraries that do not normally accompany the operating system. Such
         a contradiction means you cannot use both them and the Library together in an
         executable that you distribute.




      7. You may place library facilities that are a work based on the Library side-by-side in
         a single library together with other library facilities not covered by this License, and
         distribute such a combined library, provided that the separate distribution of the work
         based on the Library and of the other library facilities is otherwise permitted, and
         provided that you do these two things:
          ◦ Accompany the combined library with a copy of the same work based on the
            Library, uncombined with any other library facilities. This must be distributed under
            the terms of the Sections above.
          ◦ Give prominent notice with the combined library of the fact that part of it is a work
            based on the Library, and explaining where to find the accompanying uncombined
            form of the same work.
      8. You may not copy, modify, sublicense, link with, or distribute the Library except
         as expressly provided under this License. Any attempt otherwise to copy, modify,
         sublicense, link with, or distribute the Library is void, and will automatically terminate
         your rights under this License. However, parties who have received copies, or rights,
         from you under this License will not have their licenses terminated so long as such
         parties remain in full compliance.
      9. You are not required to accept this License, since you have not signed it. However,
         nothing else grants you permission to modify or distribute the Library or its derivative
         works. These actions are prohibited by law if you do not accept this License. Therefore,
         by modifying or distributing the Library (or any work based on the Library), you indicate
         your acceptance of this License to do so, and all its terms and conditions for copying,
         distributing or modifying the Library or works based on it.
     10. Each time you redistribute the Library (or any work based on the Library), the recipient
         automatically receives a license from the original licensor to copy, distribute, link with
         or modify the Library subject to these terms and conditions. You may not impose any
         further restrictions on the recipients' exercise of the rights granted herein. You are not
         responsible for enforcing compliance by third parties with this License.
     11. If, as a consequence of a court judgment or allegation of patent infringement or for any
         other reason (not limited to patent issues), conditions are imposed on you (whether by
         court order, agreement or otherwise) that contradict the conditions of this License, they
         do not excuse you from the conditions of this License. If you cannot distribute so as
         to satisfy simultaneously your obligations under this License and any other pertinent
         obligations, then as a consequence you may not distribute the Library at all. For
         example, if a patent license would not permit royalty-free redistribution of the Library by
         all those who receive copies directly or indirectly through you, then the only way you
         could satisfy both it and this License would be to refrain entirely from distribution of the
         Library.




         If any portion of this section is held invalid or unenforceable under any particular
         circumstance, the balance of the section is intended to apply, and the section as a
         whole is intended to apply in other circumstances.
         It is not the purpose of this section to induce you to infringe any patents or other
         property right claims or to contest validity of any such claims; this section has the
         sole purpose of protecting the integrity of the free software distribution system which
         is implemented by public license practices. Many people have made generous
         contributions to the wide range of software distributed through that system in reliance
         on consistent application of that system; it is up to the author/donor to decide if he or
         she is willing to distribute software through any other system and a licensee cannot
         impose that choice.
         This section is intended to make thoroughly clear what is believed to be a consequence
         of the rest of this License.
     12. If the distribution and/or use of the Library is restricted in certain countries either by
         patents or by copyrighted interfaces, the original copyright holder who places the
         Library under this License may add an explicit geographical distribution limitation
         excluding those countries, so that distribution is permitted only in or among countries
         not thus excluded. In such case, this License incorporates the limitation as if written in
         the body of this License.
     13. The Free Software Foundation may publish revised and/or new versions of the Lesser
         General Public License from time to time. Such new versions will be similar in spirit to
         the present version, but may differ in detail to address new problems or concerns.
         Each version is given a distinguishing version number. If the Library specifies a version
         number of this License which applies to it and "any later version", you have the option
         of following the terms and conditions either of that version or of any later version
         published by the Free Software Foundation. If the Library does not specify a license
         version number, you may choose any version ever published by the Free Software
         Foundation.
     14. If you wish to incorporate parts of the Library into other free programs whose
         distribution conditions are incompatible with these, write to the author to ask for
         permission. For software which is copyrighted by the Free Software Foundation, write
         to the Free Software Foundation; we sometimes make exceptions for this. Our decision
         will be guided by the two goals of preserving the free status of all derivatives of our free
         software and of promoting the sharing and reuse of software generally.
     NO WARRANTY
      • BECAUSE THE LIBRARY IS LICENSED FREE OF CHARGE, THERE IS NO
        WARRANTY FOR THE LIBRARY, TO THE EXTENT PERMITTED BY APPLICABLE
        LAW. EXCEPT WHEN OTHERWISE STATED IN WRITING THE COPYRIGHT
        HOLDERS AND/OR OTHER PARTIES PROVIDE THE LIBRARY "AS IS" WITHOUT




         WARRANTY OF ANY KIND, EITHER EXPRESSED OR IMPLIED, INCLUDING, BUT
         NOT LIMITED TO, THE IMPLIED WARRANTIES OF MERCHANTABILITY AND
         FITNESS FOR A PARTICULAR PURPOSE. THE ENTIRE RISK AS TO THE QUALITY
         AND PERFORMANCE OF THE LIBRARY IS WITH YOU. SHOULD THE LIBRARY
         PROVE DEFECTIVE, YOU ASSUME THE COST OF ALL NECESSARY SERVICING,
         REPAIR OR CORRECTION.
      • IN NO EVENT UNLESS REQUIRED BY APPLICABLE LAW OR AGREED TO
        IN WRITING WILL ANY COPYRIGHT HOLDER, OR ANY OTHER PARTY WHO
        MAY MODIFY AND/OR REDISTRIBUTE THE LIBRARY AS PERMITTED ABOVE,
        BE LIABLE TO YOU FOR DAMAGES, INCLUDING ANY GENERAL, SPECIAL,
        INCIDENTAL OR CONSEQUENTIAL DAMAGES ARISING OUT OF THE USE OR
        INABILITY TO USE THE LIBRARY (INCLUDING BUT NOT LIMITED TO LOSS OF
        DATA OR DATA BEING RENDERED INACCURATE OR LOSSES SUSTAINED BY
        YOU OR THIRD PARTIES OR A FAILURE OF THE LIBRARY TO OPERATE WITH
        ANY OTHER SOFTWARE), EVEN IF SUCH HOLDER OR OTHER PARTY HAS BEEN
        ADVISED OF THE POSSIBILITY OF SUCH DAMAGES.
\end{lstlisting}

\subsection{MariaDB}
\subsubsection*{MariaDB}

以下为 MariaDB 的许可声明原文。

\begin{lstlisting}[basicstyle=\ttfamily\footnotesize,breaklines=true]
     Synopsys modifications: None
     MariaDB is licensed under the GNU General Public License, version 2, (GPL-2.0).
     Software licensed under the GPL-2.0 is free software that comes with ABSOLUTELY NO
     WARRANTY. You may modify, redistribute, or otherwise use the GPL software under the
     terms of the GPL-2.0 (http://www.gnu.org/licenses/old-licenses/gpl-2.0.html).''
     GNU GENERAL PUBLIC LICENSE
     Version 2.1, June 1991

     Copyright (C) 1989, 1991 Free Software Foundation, Inc. 51 Franklin
     Street, Fifth Floor, Boston, MA 02110-1301, USA

     Everyone is permitted to copy and distribute verbatim copies of this
     license document, but changing it is not allowed.

     TERMS AND CONDITIONS FOR COPYING, DISTRIBUTION AND MODIFICATION
     0. This License applies to any program or other work which contains a notice placed by
     the copyright holder saying it may be distributed under the terms of this General Public
     License. The "Program", below, refers to any such program or work, and a "work based
     on the Program" means either the Program or any derivative work under copyright law:
     that is to say, a work containing the Program or a portion of it, either verbatim or with
     modifications and/or translated into another language. (Hereinafter, translation is included
     without limitation in the term "modification".) Each licensee is addressed as "you".




     Activities other than copying, distribution and modification are not covered by this License;
     they are outside its scope. The act of running the Program is not restricted, and the output
     from the Program is covered only if its contents constitute a work based on the Program
     (independent of having been made by running the Program). Whether that is true depends
     on what the Program does.
      1. You may copy and distribute verbatim copies of the Program's source code as you
         receive it, in any medium, provided that you conspicuously and appropriately publish
         on each copy an appropriate copyright notice and disclaimer of warranty; keep intact all
         the notices that refer to this License and to the absence of any warranty; and give any
         other recipients of the Program a copy of this License along with the Program.
         You may charge a fee for the physical act of transferring a copy, and you may at your
         option offer warranty protection in exchange for a fee.
      2. You may modify your copy or copies of the Program or any portion of it, thus forming a
         work based on the Program, and copy and distribute such modifications or work under
         the terms of Section 1 above, provided that you also meet all of these conditions:
          ◦ You must cause the modified files to carry prominent notices stating that you
            changed the files and the date of any change.
          ◦ You must cause any work that you distribute or publish, that in whole or in part
            contains or is derived from the Program or any part thereof, to be licensed as a
            whole at no charge to all third parties under the terms of this License.
          ◦ If the modified program normally reads commands interactively when run, you must
            cause it, when started running for such interactive use in the most ordinary way, to
            print or display an announcement including an appropriate copyright notice and a
            notice that there is no warranty (or else, saying that you provide a warranty) and
            that users may redistribute the program under these conditions, and telling the user
            how to view a copy of this License. (Exception: if the Program itself is interactive but
            does not normally print such an announcement, your work based on the Program is
            not required to print an announcement.)
             These requirements apply to the modified work as a whole. If identifiable sections
             of that work are not derived from the Program, and can be reasonably considered
             independent and separate works in themselves, then this License, and its terms, do
             not apply to those sections when you distribute them as separate works. But when
             you distribute the same sections as part of a whole which is a work based on the
             Program, the distribution of the whole must be on the terms of this License, whose
             permissions for other licensees extend to the entire whole, and thus to each and
             every part regardless of who wrote it.




             Thus, it is not the intent of this section to claim rights or contest your rights to
             work written entirely by you; rather, the intent is to exercise the right to control the
             distribution of derivative or collective works based on the Program.
             In addition, mere aggregation of another work not based on the Program with
             the Program (or with a work based on the Program) on a volume of a storage or
             distribution medium does not bring the other work under the scope of this License.
      3. You may copy and distribute the Program (or a work based on it, under Section 2) in
         object code or executable form under the terms of Sections 1 and 2 above provided
         that you also do one of the following:
         Accompany it with the complete corresponding machine-readable source code,
         which must be distributed under the terms of Sections 1 and 2 above on a medium
         customarily used for software interchange; or,
         Accompany it with a written offer, valid for at least three years, to give any third party,
         for a charge no more than your cost of physically performing source distribution, a
         complete machine-readable copy of the corresponding source code, to be distributed
         under the terms of Sections 1 and 2 above on a medium customarily used for software
         interchange; or,
         Accompany it with the information you received as to the offer to distribute
         corresponding source code. (This alternative is allowed only for noncommercial
         distribution and only if you received the program in object code or executable form with
         such an offer, in accord with Subsection b above.)
         The source code for a work means the preferred form of the work for making
         modifications to it. For an executable work, complete source code means all the source
         code for all modules it contains, plus any associated interface definition files, plus
         the scripts used to control compilation and installation of the executable. However,
         as a special exception, the source code distributed need not include anything that
         is normally distributed (in either source or binary form) with the major components
         (compiler, kernel, and so on) of the operating system on which the executable runs,
         unless that component itself accompanies the executable.
         If distribution of executable or object code is made by offering access to copy from a
         designated place, then offering equivalent access to copy the source code from the
         same place counts as distribution of the source code, even though third parties are not
         compelled to copy the source along with the object code.
      4. You may not copy, modify, sublicense, or distribute the Program except as expressly
         provided under this License. Any attempt otherwise to copy, modify, sublicense or
         distribute the Program is void, and will automatically terminate your rights under this
         License. However, parties who have received copies, or rights, from you under this
         License will not have their licenses terminated so long as such parties remain in full
         compliance.




      5. You are not required to accept this License, since you have not signed it. However,
         nothing else grants you permission to modify or distribute the Program or its derivative
         works. These actions are prohibited by law if you do not accept this License. Therefore,
         by modifying or distributing the Program (or any work based on the Program), you
         indicate your acceptance of this License to do so, and all its terms and conditions for
         copying, distributing or modifying the Program or works based on it.
      6. Each time you redistribute the Program (or any work based on the Program), the
         recipient automatically receives a license from the original licensor to copy, distribute
         or modify the Program subject to these terms and conditions. You may not impose any
         further restrictions on the recipients' exercise of the rights granted herein. You are not
         responsible for enforcing compliance by third parties to this License.
      7. If, as a consequence of a court judgment or allegation of patent infringement or for any
         other reason (not limited to patent issues), conditions are imposed on you (whether by
         court order, agreement or otherwise) that contradict the conditions of this License, they
         do not excuse you from the conditions of this License. If you cannot distribute so as
         to satisfy simultaneously your obligations under this License and any other pertinent
         obligations, then as a consequence you may not distribute the Program at all. For
         example, if a patent license would not permit royalty-free redistribution of the Program
         by all those who receive copies directly or indirectly through you, then the only way you
         could satisfy both it and this License would be to refrain entirely from distribution of the
         Program.
         If any portion of this section is held invalid or unenforceable under any particular
         circumstance, the balance of the section is intended to apply and the section as a
         whole is intended to apply in other circumstances.
         It is not the purpose of this section to induce you to infringe any patents or other
         property right claims or to contest validity of any such claims; this section has the
         sole purpose of protecting the integrity of the free software distribution system,
         which is implemented by public license practices. Many people have made generous
         contributions to the wide range of software distributed through that system in reliance
         on consistent application of that system; it is up to the author/donor to decide if he or
         she is willing to distribute software through any other system and a licensee cannot
         impose that choice.
         This section is intended to make thoroughly clear what is believed to be a consequence
         of the rest of this License.
      8. If the distribution and/or use of the Program is restricted in certain countries either
         by patents or by copyrighted interfaces, the original copyright holder who places the
         Program under this License may add an explicit geographical distribution limitation
         excluding those countries, so that distribution is permitted only in or among countries
         not thus excluded. In such case, this License incorporates the limitation as if written in
         the body of this License.




      9. The Free Software Foundation may publish revised and/or new versions of the General
         Public License from time to time. Such new versions will be similar in spirit to the
         present version, but may differ in detail to address new problems or concerns.
         Each version is given a distinguishing version number. If the Program specifies a
         version number of this License which applies to it and "any later version", you have
         the option of following the terms and conditions either of that version or of any later
         version published by the Free Software Foundation. If the Program does not specify
         a version number of this License, you may choose any version ever published by the
         Free Software Foundation.
     10. If you wish to incorporate parts of the Program into other free programs whose
         distribution conditions are different, write to the author to ask for permission. For
         software which is copyrighted by the Free Software Foundation, write to the Free
         Software Foundation; we sometimes make exceptions for this. Our decision will be
         guided by the two goals of preserving the free status of all derivatives of our free
         software and of promoting the sharing and reuse of software generally.
     NO WARRANTY
      • BECAUSE THE PROGRAM IS LICENSED FREE OF CHARGE, THERE IS NO
        WARRANTY FOR THE PROGRAM, TO THE EXTENT PERMITTED BY APPLICABLE
        LAW. EXCEPT WHEN OTHERWISE STATED IN WRITING THE COPYRIGHT
        HOLDERS AND/OR OTHER PARTIES PROVIDE THE PROGRAM "AS IS" WITHOUT
        WARRANTY OF ANY KIND, EITHER EXPRESSED OR IMPLIED, INCLUDING, BUT
        NOT LIMITED TO, THE IMPLIED WARRANTIES OF MERCHANTABILITY AND
        FITNESS FOR A PARTICULAR PURPOSE. THE ENTIRE RISK AS TO THE QUALITY
        AND PERFORMANCE OF THE PROGRAM IS WITH YOU. SHOULD THE PROGRAM
        PROVE DEFECTIVE, YOU ASSUME THE COST OF ALL NECESSARY SERVICING,
        REPAIR OR CORRECTION.
      • IN NO EVENT UNLESS REQUIRED BY APPLICABLE LAW OR AGREED TO IN
        WRITING WILL ANY COPYRIGHT HOLDER, OR ANY OTHER PARTY WHO MAY
        MODIFY AND/OR REDISTRIBUTE THE PROGRAM AS PERMITTED ABOVE,
        BE LIABLE TO YOU FOR DAMAGES, INCLUDING ANY GENERAL, SPECIAL,
        INCIDENTAL OR CONSEQUENTIAL DAMAGES ARISING OUT OF THE USE OR
        INABILITY TO USE THE PROGRAM (INCLUDING BUT NOT LIMITED TO LOSS
        OF DATA OR DATA BEING RENDERED INACCURATE OR LOSSES SUSTAINED
        BY YOU OR THIRD PARTIES OR A FAILURE OF THE PROGRAM TO OPERATE
        WITH ANY OTHER PROGRAMS), EVEN IF SUCH HOLDER OR OTHER PARTY HAS
        BEEN ADVISED OF THE POSSIBILITY OF SUCH DAMAGES.
\end{lstlisting}

\subsection{MCPP Public Domain Code}
\subsubsection*{MCPP Public Domain Code}

以下为 MCPP Public Domain Code 的许可声明原文。

\begin{lstlisting}[basicstyle=\ttfamily\footnotesize,breaklines=true]
     Copyright (c) 1998, 2002-2007 Kiyoshi Matsui <kmatsui@t3.rim.or.jp>
     All rights reserved.




     This software including the files in this directory is provided under the following license.
     Redistribution and use in source and binary forms, with or without modification, are
     permitted provided that the following conditions are met:
      1. Redistributions of source code must retain the above copyright notice, this list of
         conditions and the following disclaimer.
      2. Redistributions in binary form must reproduce the above copyright notice, this list of
         conditions and the following disclaimer in the documentation and/or other materials
         provided with the distribution.
     THIS SOFTWARE IS PROVIDED BY THE AUTHOR ''AS IS'' AND ANY EXPRESS
     OR IMPLIED WARRANTIES, INCLUDING, BUT NOT LIMITED TO, THE IMPLIED
     WARRANTIES OF MERCHANTABILITY AND FITNESS FOR A PARTICULAR PURPOSE
     ARE DISCLAIMED. IN NO EVENT SHALL THE AUTHOR BE LIABLE FOR ANY DIRECT,
     INDIRECT, INCIDENTAL, SPECIAL, EXEMPLARY, OR CONSEQUENTIAL DAMAGES
     (INCLUDING, BUT NOT LIMITED TO, PROCUREMENT OF SUBSTITUTE GOODS OR
     SERVICES; LOSS OF USE, DATA, OR PROFITS; OR BUSINESS INTERRUPTION)
     HOWEVER CAUSED AND ON ANY THEORY OF LIABILITY, WHETHER IN CONTRACT,
     STRICT LIABILITY, OR TORT (INCLUDING NEGLIGENCE OR OTHERWISE) ARISING
     IN ANY WAY OUT OF THE USE OF THIS SOFTWARE, EVEN IF ADVISED OF THE
     POSSIBILITY OF SUCH DAMAGE.
\end{lstlisting}

\subsection{NumPy*}
\subsubsection*{NumPy*}

以下为 NumPy* 的许可声明原文。

\begin{lstlisting}[basicstyle=\ttfamily\footnotesize,breaklines=true]
     Project homepage/download site: http://www.numpy.org/; https://github.com/numpy/numpy
     Project licensing notices: https://numpy.org/doc/stable/license.html
     LICENSE.txt

     Copyright (c) 2005-2017, NumPy Developers. All rights reserved.

     Redistribution and use in source and binary forms, with or without
     modification, are permitted provided that the following conditions are
     met:

     * Redistributions of source code must retain the above copyright notice,
     this list of conditions and the following disclaimer.

     * Redistributions in binary form must reproduce the above copyright
     notice, this list of conditions and the following disclaimer in the
     documentation and/or other materials provided with the distribution.

     * Neither the name of the NumPy Developers nor the names of any
     contributors may be used to endorse or promote products derived from this
     software without specific prior written permission.




     THIS SOFTWARE IS PROVIDED BY THE COPYRIGHT HOLDERS AND CONTRIBUTORS
     "AS IS" AND ANY EXPRESS OR IMPLIED WARRANTIES, INCLUDING, BUT NOT
     LIMITED TO, THE IMPLIED WARRANTIES OF MERCHANTABILITY AND FITNESS FOR
     A PARTICULAR PURPOSE ARE DISCLAIMED. IN NO EVENT SHALL THE COPYRIGHT
     OWNER OR CONTRIBUTORS BE LIABLE FOR ANY DIRECT, INDIRECT, INCIDENTAL,
     SPECIAL, EXEMPLARY, OR CONSEQUENTIAL DAMAGES (INCLUDING, BUT NOT LIMITED
     TO, PROCUREMENT OF SUBSTITUTE GOODS OR SERVICES; LOSS OF USE, DATA, OR
     PROFITS; OR BUSINESS INTERRUPTION) HOWEVER CAUSED AND ON ANY THEORY OF
     LIABILITY, WHETHER IN CONTRACT, STRICT LIABILITY, OR TORT (INCLUDING
     NEGLIGENCE OR OTHERWISE) ARISING IN ANY WAY OUT OF THE USE OF THIS
     SOFTWARE, EVEN IF ADVISED OF THE POSSIBILITY OF SUCH DAMAGE.

     The NumPy repository and source distributions bundle several libraries
     that are compatibly licensed. We list these here.

     Name: Numpydoc Files: doc/sphinxext/numpydoc/* License: 2-clause BSD For
     details, see doc/sphinxext/LICENSE.txt

     Name: scipy-sphinx-theme Files: doc/scipy-sphinx-theme/* License: 3-
     clause BSD, PSF and Apache 2.0 For details, see doc/sphinxext/LICENSE.txt

     Name: lapack-lite Files: numpy/linalg/lapack_lite/* License: 3-clause BSD
     For details, see numpy/linalg/lapack_lite/LICENSE.txt

     Name: tempita Files: tools/npy_tempita/* License: BSD derived For
     details, see tools/npy_tempita/license.txt

     Name: dragon4 Files: numpy/core/src/multiarray/dragon4.c License: One
     of a kind For license text, see numpy/core/src/multiarray/dragon4.cdoc/
     sphinxext/LICENSE.txt:

     -------------------------------------------------------------------------

     The files - numpydoc.py - docscrape.py - docscrape_sphinx.py -
     phantom_import.py have the following license: Copyright (C) 2008 Stefan
     van der Walt <stefan@mentat.za.net>, Pauli Virtanen <pav@iki.fi>

     Redistribution and use in source and binary forms, with or without
     modification, are permitted provided that the following conditions are
     met:

     1. Redistributions of source code must retain the above copyright notice,
     this list of conditions and the following disclaimer.

     2. Redistributions in binary form must reproduce the above copyright
     notice, this list of conditions and the following disclaimer in the
     documentation and/or other materials provided with the distribution.




     THIS SOFTWARE IS PROVIDED BY THE AUTHOR ``AS IS'' AND ANY EXPRESS OR
     IMPLIED WARRANTIES, INCLUDING, BUT NOT LIMITED TO, THE IMPLIED WARRANTIES
     OF MERCHANTABILITY AND FITNESS FOR A PARTICULAR PURPOSE ARE DISCLAIMED.
     IN NO EVENT SHALL THE AUTHOR BE LIABLE FOR ANY DIRECT, INDIRECT,
     INCIDENTAL, SPECIAL, EXEMPLARY, OR CONSEQUENTIAL DAMAGES (INCLUDING,
     BUT NOT LIMITED TO, PROCUREMENT OF SUBSTITUTE GOODS OR SERVICES; LOSS OF
     USE, DATA, OR PROFITS; OR BUSINESS INTERRUPTION) HOWEVER CAUSED AND ON
     ANY THEORY OF LIABILITY, WHETHER IN CONTRACT, STRICT LIABILITY, OR TORT
     (INCLUDING NEGLIGENCE OR OTHERWISE) ARISING IN ANY WAY OUT OF THE USE OF
     THIS SOFTWARE, EVEN IF ADVISED OF THE POSSIBILITY OF SUCH DAMAGE.

     -------------------------------------------------------------------------

     The files - compiler_unparse.py - comment_eater.py - traitsdoc.py have
     the following license:

     This software is OSI Certified Open Source Software. OSI Certified is
     a certification mark of the Open Source Initiative. Copyright (c) 2006,
     Enthought, Inc. All rights reserved.

     Redistribution and use in source and binary forms, with or without
     modification, are permitted provided that the following conditions are
     met:

     * Redistributions of source code must retain the above copyright notice,
     this list of conditions and the following disclaimer. * Redistributions
     in binary form must reproduce the above copyright notice, this list
     of conditions and the following disclaimer in the documentation and/or
     other materials provided with the distribution. * Neither the name of
     Enthought, Inc. nor the names of its contributors may be used to endorse
     or promote products derived from this software without specific prior
     written permission.

     THIS SOFTWARE IS PROVIDED BY THE COPYRIGHT HOLDERS AND CONTRIBUTORS
     "AS IS" AND ANY EXPRESS OR IMPLIED WARRANTIES, INCLUDING, BUT NOT
     LIMITED TO, THE IMPLIED WARRANTIES OF MERCHANTABILITY AND FITNESS FOR
     A PARTICULAR PURPOSE ARE DISCLAIMED. IN NO EVENT SHALL THE COPYRIGHT
     OWNER OR CONTRIBUTORS BE LIABLE FOR ANY DIRECT, INDIRECT, INCIDENTAL,
     SPECIAL, EXEMPLARY, OR CONSEQUENTIAL DAMAGES (INCLUDING, BUT NOT LIMITED
     TO, PROCUREMENT OF SUBSTITUTE GOODS OR SERVICES; LOSS OF USE, DATA, OR
     PROFITS; OR BUSINESS INTERRUPTION) HOWEVER CAUSED AND ON ANY THEORY OF
     LIABILITY, WHETHER IN CONTRACT, STRICT LIABILITY, OR TORT (INCLUDING
     NEGLIGENCE OR OTHERWISE) ARISING IN ANY WAY OUT OF THE USE OF THIS
     SOFTWARE, EVEN IF ADVISED OF THE POSSIBILITY OF SUCH DAMAGE.

     -------------------------------------------------------------------------




     The file - plot_directive.py originates from Matplotlib (http://
     matplotlib.sf.net/) which has the following license: Copyright (c)
     2002-2008 John D. Hunter; All Rights Reserved.

     1. This LICENSE AGREEMENT is between John D. Hunter ("JDH"), and
     the Individual or Organization ("Licensee") accessing and otherwise
     using matplotlib software in source or binary form and its associated
     documentation.

     2. Subject to the terms and conditions of this License Agreement, JDH
     hereby grants Licensee a nonexclusive, royalty-free, world-wide license
     to reproduce, analyze, test, perform and/or display publicly, prepare
     derivative works, distribute, and otherwise use matplotlib 0.98.3 alone
     or in any derivative version, provided, however, that JDH's License
     Agreement and JDH's notice of copyright, i.e., "Copyright (c) 2002-2008
     John D. Hunter; All Rights Reserved" are retained in matplotlib 0.98.3
     alone or in any derivative version prepared by Licensee.

     3. In the event Licensee prepares a derivative work that is based on or
     incorporates matplotlib 0.98.3 or any part thereof, and wants to make
     the derivative work available to others as provided herein, then Licensee
     hereby agrees to include in any such work a brief summary of the changes
     made to matplotlib 0.98.3.

     4. JDH is making matplotlib 0.98.3 available to Licensee on an "AS IS"
     basis. JDH MAKES NO REPRESENTATIONS OR WARRANTIES, EXPRESS OR IMPLIED.
     BY WAY OF EXAMPLE, BUT NOT LIMITATION, JDH MAKES NO AND DISCLAIMS
     ANY REPRESENTATION OR WARRANTY OF MERCHANTABILITY OR FITNESS FOR ANY
     PARTICULAR PURPOSE OR THAT THE USE OF MATPLOTLIB 0.98.3 WILL NOT INFRINGE
     ANY THIRD PARTY RIGHTS.

     5. JDH SHALL NOT BE LIABLE TO LICENSEE OR ANY OTHER USERS OF MATPLOTLIB
     0.98.3 FOR ANY INCIDENTAL, SPECIAL, OR CONSEQUENTIAL DAMAGES OR LOSS AS A
     RESULT OF MODIFYING, DISTRIBUTING, OR OTHERWISE USING MATPLOTLIB 0.98.3,
     OR ANY DERIVATIVE THEREOF, EVEN IF ADVISED OF THE POSSIBILITY THEREOF.

     6. This License Agreement will automatically terminate upon a material
     breach of its terms and conditions.

     7. Nothing in this License Agreement shall be deemed to create any
     relationship of agency, partnership, or joint venture between JDH and
     Licensee. This License Agreement does not grant permission to use JDH
     trademarks or trade name in a trademark sense to endorse or promote
     products or services of Licensee, or any third party.

     8. By copying, installing or otherwise using matplotlib 0.98.3, Licensee
     agrees to be bound by the terms and conditions of this License Agreement.




     numpy/linalg/lapack_lite/LICENSE.txt:

     Copyright (c) 1992-2013 The University of Tennessee and The University
     of Tennessee Research Foundation. All rights reserved. Copyright (c)
     2000-2013 The University of California Berkeley. All rights reserved.
     Copyright (c) 2006-2013 The University of Colorado Denver. All rights
     reserved.

     $COPYRIGHT$

     Additional copyrights may follow

     $HEADER$

     Redistribution and use in source and binary forms, with or without
     modification, are permitted provided that the following conditions are
     met:

     - Redistributions of source code must retain the above copyright notice,
     this list of conditions and the following disclaimer.

     - Redistributions in binary form must reproduce the above copyright
     notice, this list of conditions and the following disclaimer listed in
     this license in the documentation and/or other materials provided with
     the distribution.

     - Neither the name of the copyright holders nor the names of its
     contributors may be used to endorse or promote products derived from this
     software without specific prior written permission.

     The copyright holders provide no reassurances that the source code
     provided does not infringe any patent, copyright, or any other
     intellectual property rights of third parties. The copyright holders
     disclaim any liability to any recipient for claims brought against
     recipient by any third party for infringement of that parties
     intellectual property rights.

     THIS SOFTWARE IS PROVIDED BY THE COPYRIGHT HOLDERS AND CONTRIBUTORS
     "AS IS" AND ANY EXPRESS OR IMPLIED WARRANTIES, INCLUDING, BUT NOT
     LIMITED TO, THE IMPLIED WARRANTIES OF MERCHANTABILITY AND FITNESS FOR
     A PARTICULAR PURPOSE ARE DISCLAIMED. IN NO EVENT SHALL THE COPYRIGHT
     OWNER OR CONTRIBUTORS BE LIABLE FOR ANY DIRECT, INDIRECT, INCIDENTAL,
     SPECIAL, EXEMPLARY, OR CONSEQUENTIAL DAMAGES (INCLUDING, BUT NOT LIMITED
     TO, PROCUREMENT OF SUBSTITUTE GOODS OR SERVICES; LOSS OF USE, DATA, OR
     PROFITS; OR BUSINESS INTERRUPTION) HOWEVER CAUSED AND ON ANY THEORY OF
     LIABILITY, WHETHER IN CONTRACT, STRICT LIABILITY, OR TORT (INCLUDING
     NEGLIGENCE OR OTHERWISE) ARISING IN ANY WAY OUT OF THE USE OF THIS
     SOFTWARE, EVEN IF ADVISED OF THE POSSIBILITY OF SUCH DAMAGE.




     tools/npy_tempita/license.txt:

     License

     -------

     Copyright (c) 2008 Ian Bicking and Contributors

     Permission is hereby granted, free of charge, to any person obtaining
     a copy of this software and associated documentation files (the
     "Software"), to deal in the Software without restriction, including
     without limitation the rights to use, copy, modify, merge, publish,
     distribute, sublicense, and/or sell copies of the Software, and to
     permit persons to whom the Software is furnished to do so, subject to the
     following conditions:

     The above copyright notice and this permission notice shall be included
     in all copies or substantial portions of the Software.

     THE SOFTWARE IS PROVIDED "AS IS", WITHOUT WARRANTY OF ANY KIND,
     EXPRESS OR IMPLIED, INCLUDING BUT NOT LIMITED TO THE WARRANTIES OF
     MERCHANTABILITY, FITNESS FOR A PARTICULAR PURPOSE AND NONINFRINGEMENT. IN
     NO EVENT SHALL THE AUTHORS OR COPYRIGHT HOLDERS BE LIABLE FOR ANY CLAIM,
     DAMAGES OR OTHER LIABILITY, WHETHER IN AN ACTION OF CONTRACT, TORT OR
     OTHERWISE, ARISING FROM, OUT OF OR IN CONNECTION WITH THE SOFTWARE OR THE
     USE OR OTHER DEALINGS IN THE SOFTWARE.

     numpy/core/src/multiarray/dragon4.c:
     /*
      * Copyright (c) 2014 Ryan Juckett
      * http://www.ryanjuckett.com/
      *
      * This software is provided 'as-is', without any express or implied
      * warranty. In no event will the authors be held liable for any damages
      * arising from the use of this software.
      *
      * Permission is granted to anyone to use this software for any purpose,
      * including commercial applications, and to alter it and redistribute it
      * freely, subject to the following restrictions:
      *
      * 1. The origin of this software must not be misrepresented; you must
      not
      *    claim that you wrote the original software. If you use this
      software
      *    in a product, an acknowledgment in the product documentation would
      be
      *    appreciated but is not required.
      *
      * 2. Altered source versions must be plainly marked as such, and must
      not be




       *    misrepresented as being the original software.
       *
       * 3. This notice may not be removed or altered from any source
       *    distribution.
       */
\end{lstlisting}

\subsection{patchELF}
\subsubsection*{patchELF}

以下为 patchELF 的许可声明原文。

\begin{lstlisting}[basicstyle=\ttfamily\footnotesize,breaklines=true]
     Project homepage/download site: https://nixos.org/patchelf.html
     Project licensing notices:
     COPYING:

     See GPL-3.0 in the Standard OSS License Text on page 542."
\end{lstlisting}

\subsection{pyparsing*}
\subsubsection*{pyparsing*}

以下为 pyparsing* 的许可声明原文。

\begin{lstlisting}[basicstyle=\ttfamily\footnotesize,breaklines=true]
     Project homepage/download site: https://pypi.python.org/pypi/pyparsing
     Project licensing notices:
     # module pyparsing.py
     #
     # Copyright (c) 2003-2016 Paul T. McGuire
     #
     # Permission is hereby granted, free of charge, to any person obtaining
     # a copy of this software and associated documentation files (the
     # "Software"), to deal in the Software without restriction, including
     # without limitation the rights to use, copy, modify, merge, publish,
     # distribute, sublicense, and/or sell copies of the Software, and to
     # permit persons to whom the Software is furnished to do so, subject to
     # the following conditions:
     #
     # The above copyright notice and this permission notice shall be
     # included in all copies or substantial portions of the Software.
     #
     # THE SOFTWARE IS PROVIDED "AS IS", WITHOUT WARRANTY OF ANY KIND,
     # EXPRESS OR IMPLIED, INCLUDING BUT NOT LIMITED TO THE WARRANTIES OF
     # MERCHANTABILITY, FITNESS FOR A PARTICULAR PURPOSE AND NONINFRINGEMENT.
     # IN NO EVENT SHALL THE AUTHORS OR COPYRIGHT HOLDERS BE LIABLE FOR ANY
     # CLAIM, DAMAGES OR OTHER LIABILITY, WHETHER IN AN ACTION OF CONTRACT,
     # TORT OR OTHERWISE, ARISING FROM, OUT OF OR IN CONNECTION WITH THE
     # SOFTWARE OR THE USE OR OTHER DEALINGS IN THE SOFTWARE.
     #
\end{lstlisting}

\subsection{Python 3.6 License}
\subsubsection*{Python 3.6 License}

以下为 Python 3.6 License 的许可声明原文。

\begin{lstlisting}[basicstyle=\ttfamily\footnotesize,breaklines=true]
     PYTHON SOFTWARE FOUNDATION LICENSE VERSION 2




     --------------------------------------------

      1. This LICENSE AGREEMENT is between the Python Software Foundation
         ("PSF"), and the Individual or Organization ("Licensee") accessing and
         otherwise using this software ("Python") in source or binary form and
         its associated documentation.

      2. Subject to the terms and conditions of this License Agreement, PSF
         hereby grants Licensee a nonexclusive, royalty-free, world-wide
         license to reproduce, analyze, test, perform and/or display publicly,
         prepare derivative works, distribute, and otherwise use Python alone
         or in any derivative version, provided, however, that PSF's License
         Agreement and PSF's notice of copyright, i.e., "Copyright (c) 2001,
         2002, 2003, 2004, 2005, 2006, 2007, 2008, 2009, 2010, 2011, 2012,
         2013, 2014 2015, 2016 Python Software Foundation; All Rights Reserved"
         are retained in Python alone or in any derivative version prepared by
         Licensee.

      3. In the event Licensee prepares a derivative work that is based on
         or incorporates Python or any part thereof, and wants to make the
         derivative work available to others as provided herein, then Licensee
         hereby agrees to include in any such work a brief summary of the
         changes made to Python.

      4. PSF is making Python available to Licensee on an "AS IS" basis.
         PSF MAKES NO REPRESENTATIONS OR WARRANTIES, EXPRESS OR IMPLIED. BY
         WAY OF EXAMPLE, BUT NOT LIMITATION, PSF MAKES NO AND DISCLAIMS ANY
         REPRESENTATION OR WARRANTY OF MERCHANTABILITY OR FITNESS FOR ANY
         PARTICULAR PURPOSE OR THAT THE USE OF PYTHON WILL NOT INFRINGE ANY
         THIRD PARTY RIGHTS.

      5. PSF SHALL NOT BE LIABLE TO LICENSEE OR ANY OTHER USERS OF PYTHON
         FOR ANY INCIDENTAL, SPECIAL, OR CONSEQUENTIAL DAMAGES OR LOSS AS A
         RESULT OF MODIFYING, DISTRIBUTING, OR OTHERWISE USING PYTHON, OR ANY
         DERIVATIVE THEREOF, EVEN IF ADVISED OF THE POSSIBILITY THEREOF.

      6. This License Agreement will automatically terminate upon a material
         breach of its terms and conditions.

      7. Nothing in this License Agreement shall be deemed to create any
         relationship of agency, partnership, or joint venture between PSF and
         Licensee. This License Agreement does not grant permission to use PSF
         trademarks or trade name in a trademark sense to endorse or promote
         products or services of Licensee, or any third party.

      8. By copying, installing or otherwise using Python, Licensee agrees to
         be bound by the terms and conditions of this License Agreement.
\end{lstlisting}

\subsection{scikit-learn}
\subsubsection*{scikit-learn}

以下为 scikit-learn 的许可声明原文。

\begin{lstlisting}[basicstyle=\ttfamily\footnotesize,breaklines=true]
     Project homepage/download site: http://scikit-learn.org/
     Project licensing notices:
     New BSD License

     Copyright (c) 2007-2018 The scikit-learn developers. All rights reserved.
     Redistribution and use in source and binary forms, with or without
     modification, are permitted provided that the following conditions are
     met:

     a. Redistributions of source code must retain the above copyright notice,
     this list of conditions and the following disclaimer.

     b. Redistributions in binary form must reproduce the above copyright
     notice, this list of conditions and the following disclaimer in the
     documentation and/or other materials provided with the distribution.

     c. Neither the name of the Scikit-learn Developers nor the names of its
     contributors may be used to endorse or promote products derived from this
     software without specific prior written permission.

     THIS SOFTWARE IS PROVIDED BY THE COPYRIGHT HOLDERS AND CONTRIBUTORS "AS
     IS" AND ANY EXPRESS OR IMPLIED WARRANTIES, INCLUDING, BUT NOT LIMITED TO,
     THE IMPLIED WARRANTIES OF MERCHANTABILITY AND FITNESS FOR A PARTICULAR
     PURPOSE ARE DISCLAIMED. IN NO EVENT SHALL THE REGENTS OR CONTRIBUTORS
     BE LIABLE FOR ANY DIRECT, INDIRECT, INCIDENTAL, SPECIAL, EXEMPLARY, OR
     CONSEQUENTIAL DAMAGES (INCLUDING, BUT NOT LIMITED TO, PROCUREMENT OF
     SUBSTITUTE GOODS OR SERVICES; LOSS OF USE, DATA, OR PROFITS; OR BUSINESS
     INTERRUPTION) HOWEVER CAUSED AND ON ANY THEORY OF LIABILITY, WHETHER IN
     CONTRACT, STRICT LIABILITY, OR TORT (INCLUDING NEGLIGENCE OR OTHERWISE)
     ARISING IN ANY WAY OUT OF THE USE OF THIS SOFTWARE, EVEN IF ADVISED OF
     THE POSSIBILITY OF SUCH DAMAGE.
\end{lstlisting}

\subsection{SciPy*}
\subsubsection*{SciPy*}

以下为 SciPy* 的许可声明原文。

\begin{lstlisting}[basicstyle=\ttfamily\footnotesize,breaklines=true]
     Project homepage/download site: https://www.SciPy.org; https://github.com/SciPy/SciPy
     Project licensing notices: https://github.com/SciPy/SciPy/blob/master/LICENSE.txt
     Copyright (c) 2001, 2002 Enthought, Inc. All rights reserved.
     Copyright (c) 2003-2017 SciPy Developers. All rights reserved.
     Redistribution and use in source and binary forms, with or without
      modification, are permitted provided that the following conditions are
      met:




     a. Redistributions of source code must retain the above copyright notice,
      this list of conditions and the following disclaimer.
     b. Redistributions in binary form must reproduce the above copyright
      notice, this list of conditions and the following disclaimer in the
      documentation and/or other materials provided with the distribution.
     c. Neither the name of Enthought nor the names of the SciPy Developers
      may be used to endorse or promote products derived from this software
      without specific prior written permission. THIS SOFTWARE IS PROVIDED
      BY THE COPYRIGHT HOLDERS AND CONTRIBUTORS "AS IS" AND ANY EXPRESS
      OR IMPLIED WARRANTIES, INCLUDING, BUT NOT LIMITED TO, THE IMPLIED
      WARRANTIES OF MERCHANTABILITY AND FITNESS FOR A PARTICULAR PURPOSE ARE
      DISCLAIMED. IN NO EVENT SHALL THE COPYRIGHT HOLDERS OR CONTRIBUTORS
      BE LIABLE FOR ANY DIRECT, INDIRECT, INCIDENTAL, SPECIAL, EXEMPLARY, OR
      CONSEQUENTIAL DAMAGES (INCLUDING, BUT NOT LIMITED TO, PROCUREMENT OF
      SUBSTITUTE GOODS OR SERVICES; LOSS OF USE, DATA, OR PROFITS; OR BUSINESS
      INTERRUPTION) HOWEVER CAUSED AND ON ANY THEORY OF LIABILITY, WHETHER IN
      CONTRACT, STRICT LIABILITY, OR TORT (INCLUDING NEGLIGENCE OR OTHERWISE)
      ARISING IN ANY WAY OUT OF THE USE OF THIS SOFTWARE, EVEN IF ADVISED OF
      THE POSSIBILITY OF SUCH DAMAGE.
     SciPy bundles a number of libraries that are compatibly licensed. We
      list these here.
     Name: Numpydoc
     Files: doc/sphinxext/numpydoc/*
     License: 2-clause BSD
       For details, see doc/sphinxext/LICENSE.txt
     Name: scipy-sphinx-theme
     Files: doc/scipy-sphinx-theme/*
     License: 3-clause BSD, PSF and Apache 2.0
       For details, see doc/sphinxext/LICENSE.txt
     Name: Six
     Files: scipy/_lib/six.py
     License: MIT
       For details, see the header inside scipy/_lib/six.py
     Name: Decorator
     Files: scipy/_lib/decorator.py
     License: 2-clause BSD
       For details, see the header inside scipy/_lib/decorator.py
     Name: ID
     Files: scipy/linalg/src/id_dist/*
     License: 3-clause BSD
       For details, see scipy/linalg/src/id_dist/doc/doc.tex
     Name: L-BFGS-B
     Files: scipy/optimize/lbfgsb/*
     License: BSD license
       For details, see scipy/optimize/lbfgsb/README
     Name: SuperLU
     Files: scipy/sparse/linalg/dsolve/SuperLU/*
     License: 3-clause BSD
       For details, see scipy/sparse/linalg/dsolve/SuperLU/License.txt
     Name: ARPACK
     Files: scipy/sparse/linalg/eigen/arpack/ARPACK/*
     License: 3-clause BSD
       For details, see scipy/sparse/linalg/eigen/arpack/ARPACK/COPYING




     Name: Qhull
     Files: scipy/spatial/qhull/*
     License: Qhull license (BSD-like)
       For details, see scipy/spatial/qhull/COPYING.txt
     Name: Cephes
     Files: scipy/special/cephes/*
     License: 3-clause BSD
       Distributed under 3-clause BSD license with permission from the author,
       see https://lists.debian.org/debian-legal/2004/12/msg00295.html
       Cephes Math Library Release 2.8: June, 2000
       Copyright 1984, 1995, 2000 by Stephen L. Moshier
       This software is derived from the Cephes Math Library and is
       incorporated herein by permission of the author.
       All rights reserved.
       Redistribution and use in source and binary forms, with or without
      modification, are permitted provided that the following conditions are
      met:
     * Redistributions of source code must retain the above copyright notice,
      this list of conditions and the following disclaimer.
     * Redistributions in binary form must reproduce the above copyright
      notice, this list of conditions and the following disclaimer in the
      documentation and/or other materials provided with the distribution.
     * Neither the name of the <organization> nor the names of its
      contributors may be used to endorse or promote products derived from
      this software without specific prior written permission.
     THIS SOFTWARE IS PROVIDED BY THE COPYRIGHT HOLDERS AND CONTRIBUTORS
      "AS IS" AND ANY EXPRESS OR IMPLIED WARRANTIES, INCLUDING, BUT NOT
      LIMITED TO, THE IMPLIED WARRANTIES OF MERCHANTABILITY AND FITNESS FOR A
      PARTICULAR PURPOSE ARE DISCLAIMED. IN NO EVENT SHALL <COPYRIGHT HOLDER>
      BE LIABLE FOR ANY DIRECT, INDIRECT, INCIDENTAL, SPECIAL, EXEMPLARY, OR
      CONSEQUENTIAL DAMAGES (INCLUDING, BUT NOT LIMITED TO, PROCUREMENT OF
      SUBSTITUTE GOODS OR SERVICES; LOSS OF USE, DATA, OR PROFITS; OR BUSINESS
      INTERRUPTION) HOWEVER CAUSED AND ON ANY THEORY OF LIABILITY, WHETHER IN
      CONTRACT, STRICT LIABILITY, OR TORT (INCLUDING NEGLIGENCE OR OTHERWISE)
      ARISING IN ANY WAY OUT OF THE USE OF THIS SOFTWARE, EVEN IF ADVISED OF
      THE POSSIBILITY OF SUCH DAMAGE.
     Name: Faddeeva
     Files: scipy/special/Faddeeva.*
     License: MIT
       Copyright (c) 2012 Massachusetts Institute of Technology
       Permission is hereby granted, free of charge, to any person obtaining a
     copy of this software and associated documentation files (the
     "Software"), to deal in the Software without restriction, including
     without limitation the rights to use, copy, modify, merge, publish,
     distribute, sublicense, and/or sell copies of the Software, and to permit
     persons to whom the Software is furnished to do so, subject to the
     following conditions:
     The above copyright notice and this permission notice shall be included
     in all copies or substantial portions of the Software.
     THE SOFTWARE IS PROVIDED "AS IS", WITHOUT WARRANTY OF ANY KIND, EXPRESS
     OR IMPLIED, INCLUDING BUT NOT LIMITED TO THE WARRANTIES OF
     MERCHANTABILITY, FITNESS FOR A PARTICULAR PURPOSE AND NONINFRINGEMENT. IN
     NO EVENT SHALL THE AUTHORS OR COPYRIGHT HOLDERS BE LIABLE FOR ANY CLAIM,




     DAMAGES OR OTHER LIABILITY, WHETHER IN AN ACTION OF CONTRACT, TORT OR
     OTHERWISE, ARISING FROM, OUT OF OR IN CONNECTION WITH THE SOFTWARE OR THE
     USE OR OTHER DEALINGS IN THE SOFTWARE.
     doc/sphinxext/LICENSE.txt:
     -------------------------------------------------------------------------
     ------
     The files - numpydoc.py - docscrape.py - docscrape_sphinx.py -
     phantom_import.py have the following license:
     Copyright (C) 2008 Stefan van der Walt <stefan@mentat.za.net>, Pauli
     Virtanen <pav@iki.fi>
     Redistribution and use in source and binary forms, with or without
      modification, are permitted provided that the following conditions are
      met:
     1. Redistributions of source code must retain the above copyright notice,
      this list of conditions and the following disclaimer.
     2. Redistributions in binary form must reproduce the above copyright
      notice, this list of conditions and the following disclaimer in the
      documentation and/or other materials provided with the distribution.
     THIS SOFTWARE IS PROVIDED BY THE AUTHOR ``AS IS'' AND ANY EXPRESS
      OR IMPLIED WARRANTIES, INCLUDING, BUT NOT LIMITED TO, THE IMPLIED
      WARRANTIES OF MERCHANTABILITY AND FITNESS FOR A PARTICULAR PURPOSE
      ARE DISCLAIMED. IN NO EVENT SHALL THE AUTHOR BE LIABLE FOR ANY DIRECT,
      INDIRECT, INCIDENTAL, SPECIAL, EXEMPLARY, OR CONSEQUENTIAL DAMAGES
      (INCLUDING, BUT NOT LIMITED TO, PROCUREMENT OF SUBSTITUTE GOODS OR
      SERVICES; LOSS OF USE, DATA, OR PROFITS; OR BUSINESS INTERRUPTION)
      HOWEVER CAUSED AND ON ANY THEORY OF LIABILITY, WHETHER IN CONTRACT,
      STRICT LIABILITY, OR TORT (INCLUDING NEGLIGENCE OR OTHERWISE) ARISING
      IN ANY WAY OUT OF THE USE OF THIS SOFTWARE, EVEN IF ADVISED OF THE
      POSSIBILITY OF SUCH DAMAGE.
     -------------------------------------------------------------------------
     ------
     The files - compiler_unparse.py - comment_eater.py - traitsdoc.py have
      the following license:
     This software is OSI Certified Open Source Software. OSI Certified is a
      certification mark of the Open Source Initiative.
     Copyright (c) 2006, Enthought, Inc. All rights reserved.
     Redistribution and use in source and binary forms, with or without
      modification, are permitted provided that the following conditions are
      met:
      * Redistributions of source code must retain the above copyright notice,
      this list of conditions and the following disclaimer. * Redistributions
      in binary form must reproduce the above copyright notice, this list of
      conditions and the following disclaimer in the documentation and/or
      other materials provided with the distribution. * Neither the name of
      Enthought, Inc. nor the names of its contributors may be used to endorse
      or promote products derived from this software without specific prior
      written permission.
     THIS SOFTWARE IS PROVIDED BY THE COPYRIGHT HOLDERS AND CONTRIBUTORS
      "AS IS" AND ANY EXPRESS OR IMPLIED WARRANTIES, INCLUDING, BUT NOT
      LIMITED TO, THE IMPLIED WARRANTIES OF MERCHANTABILITY AND FITNESS FOR
      A PARTICULAR PURPOSE ARE DISCLAIMED. IN NO EVENT SHALL THE COPYRIGHT
      OWNER OR CONTRIBUTORS BE LIABLE FOR ANY DIRECT, INDIRECT, INCIDENTAL,
      SPECIAL, EXEMPLARY, OR CONSEQUENTIAL DAMAGES (INCLUDING, BUT NOT LIMITED




      TO, PROCUREMENT OF SUBSTITUTE GOODS OR SERVICES; LOSS OF USE, DATA, OR
      PROFITS; OR BUSINESS INTERRUPTION) HOWEVER CAUSED AND ON ANY THEORY OF
      LIABILITY, WHETHER IN CONTRACT, STRICT LIABILITY, OR TORT (INCLUDING
      NEGLIGENCE OR OTHERWISE) ARISING IN ANY WAY OUT OF THE USE OF THIS
      SOFTWARE, EVEN IF ADVISED OF THE POSSIBILITY OF SUCH DAMAGE.
     -------------------------------------------------------------------------
     ------
     The file - plot_directive.py originates from Matplotlib
      (http://matplotlib.sf.net/) which has the following license:
     Copyright (c) 2002-2008 John D. Hunter; All Rights Reserved.
     1. This LICENSE AGREEMENT is between John D. Hunter ("JDH"), and
      the Individual or Organization ("Licensee") accessing and otherwise
      using matplotlib software in source or binary form and its associated
      documentation.
     2. Subject to the terms and conditions of this License Agreement, JDH
      hereby grants Licensee a nonexclusive, royalty-free, world-wide license
      to reproduce, analyze, test, perform and/or display publicly, prepare
      derivative works, distribute, and otherwise use matplotlib 0.98.3 alone
      or in any derivative version, provided, however, that JDH's License
      Agreement and JDH's notice of copyright, i.e., "Copyright (c) 2002-2008
      John D. Hunter; All Rights Reserved" are retained in matplotlib 0.98.3
      alone or in any derivative version prepared by Licensee.
     3. In the event Licensee prepares a derivative work that is based on
      or incorporates matplotlib 0.98.3 or any part thereof, and wants to
      make the derivative work available to others as provided herein, then
      Licensee hereby agrees to include in any such work a brief summary of
      the changes made to matplotlib 0.98.3.
     4. JDH is making matplotlib 0.98.3 available to Licensee on an "AS IS"
      basis. JDH MAKES NO REPRESENTATIONS OR WARRANTIES, EXPRESS OR IMPLIED.
      BY WAY OF EXAMPLE, BUT NOT LIMITATION, JDH MAKES NO AND DISCLAIMS
      ANY REPRESENTATION OR WARRANTY OF MERCHANTABILITY OR FITNESS FOR
      ANY PARTICULAR PURPOSE OR THAT THE USE OF MATPLOTLIB 0.98.3 WILL NOT
      INFRINGE ANY THIRD PARTY RIGHTS.
     5. JDH SHALL NOT BE LIABLE TO LICENSEE OR ANY OTHER USERS OF MATPLOTLIB
      0.98.3 FOR ANY INCIDENTAL, SPECIAL, OR CONSEQUENTIAL DAMAGES OR LOSS
      AS A RESULT OF MODIFYING, DISTRIBUTING, OR OTHERWISE USING MATPLOTLIB
      0.98.3, OR ANY DERIVATIVE THEREOF, EVEN IF ADVISED OF THE POSSIBILITY
      THEREOF.
     6. This License Agreement will automatically terminate upon a material
      breach of its terms and conditions.
     7. Nothing in this License Agreement shall be deemed to create any
      relationship of agency, partnership, or joint venture between JDH and
      Licensee. This License Agreement does not grant permission to use JDH
      trademarks or trade name in a trademark sense to endorse or promote
      products or services of Licensee, or any third party.
     8. By copying, installing or otherwise using matplotlib 0.98.3, Licensee
      agrees to be bound by the terms and conditions of this License
      Agreement.
     numpy/linalg/lapack_lite/LICENSE.txt:
     Copyright (c) 1992-2013 The University of Tennessee and The University
      of Tennessee Research Foundation. All rights reserved. Copyright (c)
      2000-2013 The University of California Berkeley. All rights reserved.




      Copyright (c) 2006-2013 The University of Colorado Denver. All rights
      reserved.
     $COPYRIGHT$
     Additional copyrights may follow
     $HEADER$
     Redistribution and use in source and binary forms, with or without
      modification, are permitted provided that the following conditions are
      met:
     - Redistributions of source code must retain the above copyright notice,
      this list of conditions and the following disclaimer.
     - Redistributions in binary form must reproduce the above copyright
      notice, this list of conditions and the following disclaimer listed in
      this license in the documentation and/or other materials provided with
      the distribution.
     - Neither the name of the copyright holders nor the names of its
      contributors may be used to endorse or promote products derived from
      this software without specific prior written permission.
     The copyright holders provide no reassurances that the source code
      provided does not infringe any patent, copyright, or any other
      intellectual property rights of third parties. The copyright holders
      disclaim any liability to any recipient for claims brought against
      recipient by any third party for infringement of that parties
      intellectual property rights.
     THIS SOFTWARE IS PROVIDED BY THE COPYRIGHT HOLDERS AND CONTRIBUTORS
      "AS IS" AND ANY EXPRESS OR IMPLIED WARRANTIES, INCLUDING, BUT NOT
      LIMITED TO, THE IMPLIED WARRANTIES OF MERCHANTABILITY AND FITNESS FOR
      A PARTICULAR PURPOSE ARE DISCLAIMED. IN NO EVENT SHALL THE COPYRIGHT
      OWNER OR CONTRIBUTORS BE LIABLE FOR ANY DIRECT, INDIRECT, INCIDENTAL,
      SPECIAL, EXEMPLARY, OR CONSEQUENTIAL DAMAGES (INCLUDING, BUT NOT LIMITED
      TO, PROCUREMENT OF SUBSTITUTE GOODS OR SERVICES; LOSS OF USE, DATA, OR
      PROFITS; OR BUSINESS INTERRUPTION) HOWEVER CAUSED AND ON ANY THEORY OF
      LIABILITY, WHETHER IN CONTRACT, STRICT LIABILITY, OR TORT (INCLUDING
      NEGLIGENCE OR OTHERWISE) ARISING IN ANY WAY OUT OF THE USE OF THIS
      SOFTWARE, EVEN IF ADVISED OF THE POSSIBILITY OF SUCH DAMAGE.
     tools/npy_tempita/license.txt:
     License
     -------
     Copyright (c) 2008 Ian Bicking and Contributors
     Permission is hereby granted, free of charge, to any person obtaining
      a copy of this software and associated documentation files (the
      "Software"), to deal in the Software without restriction, including
      without limitation the rights to use, copy, modify, merge, publish,
      distribute, sublicense, and/or sell copies of the Software, and to
      permit
     persons to whom the Software is furnished to do so, subject to the
     following conditions:
     The above copyright notice and this permission notice shall be included
      in all copies or substantial portions of the Software.
     THE SOFTWARE IS PROVIDED "AS IS", WITHOUT WARRANTY OF ANY KIND,
      EXPRESS OR IMPLIED, INCLUDING BUT NOT LIMITED TO THE WARRANTIES OF
      MERCHANTABILITY, FITNESS FOR A PARTICULAR PURPOSE AND NONINFRINGEMENT.
      IN NO EVENT SHALL THE AUTHORS OR COPYRIGHT HOLDERS BE LIABLE FOR ANY
      CLAIM, DAMAGES OR OTHER LIABILITY, WHETHER IN AN ACTION OF CONTRACT,




       TORT OR OTHERWISE, ARISING FROM, OUT OF OR IN CONNECTION WITH THE
       SOFTWARE OR THE USE OR OTHER DEALINGS IN THE SOFTWARE.
     scipy/_lib/six.py
     """Utilities for writing code that runs on Python 2 and 3"""
     # Copyright (c) 2010-2012 Benjamin Peterson
     #
     # Permission is hereby granted, free of charge, to any person obtaining a
       copy of
     # this software and associated documentation files (the "Software"), to
       deal in
     # the Software without restriction, including without limitation the
     rights to
     # use, copy, modify, merge, publish, distribute, sublicense, and/or sell
     copies of
     # the Software, and to permit persons to whom the Software is furnished
     to do so,
     # subject to the following conditions:
     #
     # The above copyright notice and this permission notice shall be included
     in all
     # copies or substantial portions of the Software.
     #
     # THE SOFTWARE IS PROVIDED "AS IS", WITHOUT WARRANTY OF ANY KIND, EXPRESS
     OR
     # IMPLIED, INCLUDING BUT NOT LIMITED TO THE WARRANTIES OF
     MERCHANTABILITY, FITNESS
     # FOR A PARTICULAR PURPOSE AND NONINFRINGEMENT. IN NO EVENT SHALL THE
     AUTHORS OR
     # COPYRIGHT HOLDERS BE LIABLE FOR ANY CLAIM, DAMAGES OR OTHER LIABILITY,
     WHETHER
     # IN AN ACTION OF CONTRACT, TORT OR OTHERWISE, ARISING FROM, OUT OF OR IN
     # CONNECTION WITH THE SOFTWARE OR THE USE OR OTHER DEALINGS IN THE
     SOFTWARE.
     scipy/_lib/decorator.py
     # #########################      LICENSE     ############################
     #
     # Copyright (c) 2005-2015, Michele Simionato
     # All rights reserved.
     # Redistribution and use in source and binary forms, with or without
     # modification, are permitted provided that the following conditions are
     # met:
     #    Redistributions of source code must retain the above copyright
     #    notice, this list of conditions and the following disclaimer.
     #    Redistributions in bytecode form must reproduce the above copyright
     #    notice, this list of conditions and the following disclaimer in
     #    the documentation and/or other materials provided with the
     #    distribution.
     # THIS SOFTWARE IS PROVIDED BY THE COPYRIGHT HOLDERS AND CONTRIBUTORS
     # "AS IS" AND ANY EXPRESS OR IMPLIED WARRANTIES, INCLUDING, BUT NOT
     # LIMITED TO, THE IMPLIED WARRANTIES OF MERCHANTABILITY AND FITNESS FOR
     # A PARTICULAR PURPOSE ARE DISCLAIMED. IN NO EVENT SHALL THE COPYRIGHT
     # HOLDERS OR CONTRIBUTORS BE LIABLE FOR ANY DIRECT, INDIRECT,
     # INCIDENTAL, SPECIAL, EXEMPLARY, OR CONSEQUENTIAL DAMAGES (INCLUDING,




     # BUT NOT LIMITED TO, PROCUREMENT OF SUBSTITUTE GOODS OR SERVICES; LOSS
     # OF USE, DATA, OR PROFITS; OR BUSINESS INTERRUPTION) HOWEVER CAUSED AND
     # ON ANY THEORY OF LIABILITY, WHETHER IN CONTRACT, STRICT LIABILITY, OR
     # TORT (INCLUDING NEGLIGENCE OR OTHERWISE) ARISING IN ANY WAY OUT OF THE
     # USE OF THIS SOFTWARE, EVEN IF ADVISED OF THE POSSIBILITY OF SUCH
     # DAMAGE.
     scipy/linalg/src/id_dist/doc/doc.tex
     The present document and all of the software
     in the accompanying distribution (which is contained in the directory
     {\tt id\_dist} and its subdirectories, or in the file
     {\tt id\_dist.tar.gz})\, is
     \bigskip
     Copyright \copyright\ 2014 by P.-G. Martinsson, V. Rokhlin,
     Y. Shkolnisky, and M. Tygert.
     \bigskip
     All rights reserved.
     \bigskip
     Redistribution and use in source and binary forms, with or without
     modification, are permitted provided that the following conditions are
     met:
     \begin{enumerate}
     \item Redistributions of source code must retain the above copyright
     notice, this list of conditions, and the following disclaimer.
     \item Redistributions in binary form must reproduce the above copyright
     notice, this list of conditions, and the following disclaimer in the
     documentation and/or other materials provided with the distribution.
     \item None of the names of the copyright holders may be used to endorse
     or promote products derived from this software without specific prior
     written permission.
     \end{enumerate}
     \bigskip
     THIS SOFTWARE IS PROVIDED BY THE COPYRIGHT HOLDERS ``AS IS'' AND ANY
     EXPRESS OR IMPLIED WARRANTIES, INCLUDING, BUT NOT LIMITED TO, THE
     IMPLIED WARRANTIES OF MERCHANTABILITY AND FITNESS FOR A PARTICULAR
     PURPOSE ARE DISCLAIMED. IN NO EVENT SHALL THE COPYRIGHT OWNERS BE
     LIABLE FOR ANY DIRECT, INDIRECT, INCIDENTAL, SPECIAL, EXEMPLARY, OR
     CONSEQUENTIAL DAMAGES (INCLUDING, BUT NOT LIMITED TO,
     PROCUREMENT

      OF
     SUBSTITUTE GOODS OR SERVICES; LOSS OF USE, DATA, OR PROFITS; OR
     BUSINESS INTERRUPTION) HOWEVER CAUSED AND ON ANY THEORY OF LIABILITY,
     WHETHER IN CONTRACT, STRICT LIABILITY, OR TORT (INCLUDING NEGLIGENCE
      OR
     OTHERWISE) ARISING IN ANY WAY OUT OF THE USE OF THIS SOFTWARE, EVEN IF
     ADVISED OF THE POSSIBILITY OF SUCH DAMAGE.
     scipy/optimize/lbfgsb/README
     License of L-BFGS-B (Fortran code)
     ==================================
     The version included here (in lbfgsb.f) is 3.0 (released April 25, 2011).
     It was written by Ciyou Zhu, Richard Byrd, and Jorge Nocedal
     <nocedal@ece.nwu.edu>. It carries the following condition for use:
     """




     This software is freely available, but we expect that all publications
     describing work using this software, or all commercial products using it,
     quote at least one of the references given below. This software is
     released under the BSD License.
     References
     * R. H. Byrd, P. Lu and J. Nocedal. A Limited Memory Algorithm for Bound
     Constrained Optimization, (1995), SIAM Journal on Scientific and
     Statistical Computing, 16, 5, pp. 1190-1208.
     * C. Zhu, R. H. Byrd and J. Nocedal. L-BFGS-B: Algorithm 778: L-BFGS-B,
     FORTRAN routines for large scale bound constrained optimization (1997),
     ACM Transactions on Mathematical Software, 23, 4, pp. 550 - 560.
     * J.L. Morales and J. Nocedal. L-BFGS-B: Remark on Algorithm 778:
     L-BFGS-B, FORTRAN routines for large scale bound constrained optimization
     (2011), ACM Transactions on Mathematical Software, 38, 1.
     """
     License for the Python wrapper
     ==============================
     Copyright (c) 2004 David M. Cooke <cookedm@physics.mcmaster.ca>
     Permission is hereby granted, free of charge, to any person obtaining a
     copy of this software and associated documentation files (the
     "Software"), to deal in the Software without restriction, including
     without limitation the rights to use, copy, modify, merge, publish,
     distribute, sublicense, and/or sell copies of the Software, and to permit
     persons to whom the Software is furnished to do so, subject to the
     following conditions:
     The above copyright notice and this permission notice shall be included
     in all copies or substantial portions of the Software.
     THE SOFTWARE IS PROVIDED "AS IS", WITHOUT WARRANTY OF ANY KIND, EXPRESS
     OR IMPLIED, INCLUDING BUT NOT LIMITED TO THE WARRANTIES OF
     MERCHANTABILITY, FITNESS FOR A PARTICULAR PURPOSE AND NONINFRINGEMENT. IN
     NO EVENT SHALL THE AUTHORS OR COPYRIGHT HOLDERS BE LIABLE FOR ANY CLAIM,
     DAMAGES OR OTHER LIABILITY, WHETHER IN AN ACTION OF CONTRACT, TORT OR
     OTHERWISE, ARISING FROM, OUT OF OR IN CONNECTION WITH THE SOFTWARE OR THE
     USE OR OTHER DEALINGS IN THE SOFTWARE.
     scipy/sparse/linalg/dsolve/SuperLU/License.txt
     Copyright (c) 2003, The Regents of the University of California, through
     Lawrence Berkeley National Laboratory (subject to receipt of any required
     approvals from U.S. Dept. of Energy)
     All rights reserved.
     Redistribution and use in source and binary forms, with or without
     modification, are permitted provided that the following conditions are
     met:
     (1) Redistributions of source code must retain the above copyright
      notice, this list of conditions and the following disclaimer.
     (2) Redistributions in binary form must reproduce the above copyright
      notice, this list of conditions and the following disclaimer in the
      documentation and/or other materials provided with the distribution.
     (3) Neither the name of Lawrence Berkeley National Laboratory, U.S. Dept.
      of Energy nor the names of its contributors may be used to endorse
      or promote products derived from this software without specific prior
      written permission.
     THIS SOFTWARE IS PROVIDED BY THE COPYRIGHT HOLDERS AND CONTRIBUTORS
      "AS IS" AND ANY EXPRESS OR IMPLIED WARRANTIES, INCLUDING, BUT NOT




      LIMITED TO, THE IMPLIED WARRANTIES OF MERCHANTABILITY AND FITNESS FOR
      A PARTICULAR PURPOSE ARE DISCLAIMED. IN NO EVENT SHALL THE COPYRIGHT
      OWNER OR CONTRIBUTORS BE LIABLE FOR ANY DIRECT, INDIRECT, INCIDENTAL,
      SPECIAL, EXEMPLARY, OR CONSEQUENTIAL DAMAGES (INCLUDING, BUT NOT LIMITED
      TO, PROCUREMENT OF SUBSTITUTE GOODS OR SERVICES; LOSS OF USE, DATA, OR
      PROFITS; OR BUSINESS INTERRUPTION) HOWEVER CAUSED AND ON ANY THEORY OF
      LIABILITY, WHETHER IN CONTRACT, STRICT LIABILITY, OR TORT (INCLUDING
      NEGLIGENCE OR OTHERWISE) ARISING IN ANY WAY OUT OF THE USE OF THIS
      SOFTWARE, EVEN IF ADVISED OF THE POSSIBILITY OF SUCH DAMAGE.
     scipy/sparse/linalg/eigen/arpack/ARPACK/COPYING
     BSD Software License
     Pertains to ARPACK and P_ARPACK
     Copyright (c) 1996-2008 Rice University.
     Developed by D.C. Sorensen, R.B. Lehoucq, C. Yang, and K. Maschhoff.
     All rights reserved.
     Arpack has been renamed to arpack-ng.
     Copyright (c) 2001-2011 - Scilab Enterprises
     Updated by Allan Cornet, Sylvestre Ledru.
     Copyright (c) 2010 - Jordi Gutiérrez Hermoso (Octave patch)
     Copyright (c) 2007 - Sébastien Fabbro (gentoo patch)
     Redistribution and use in source and binary forms, with or without
      modification, are permitted provided that the following conditions are
      met:
     - Redistributions of source code must retain the above copyright notice,
      this list of conditions and the following disclaimer.
     - Redistributions in binary form must reproduce the above copyright
      notice, this list of conditions and the following disclaimer listed in
      this license in the documentation and/or other materials provided with
      the distribution.
     - Neither the name of the copyright holders nor the names of its
      contributors may be used to endorse or promote products derived from
      this software without specific prior written permission.
     THIS SOFTWARE IS PROVIDED BY THE COPYRIGHT HOLDERS AND CONTRIBUTORS
      "AS IS" AND ANY EXPRESS OR IMPLIED WARRANTIES, INCLUDING, BUT NOT
      LIMITED TO, THE IMPLIED WARRANTIES OF MERCHANTABILITY AND FITNESS FOR
      A PARTICULAR PURPOSE ARE DISCLAIMED. IN NO EVENT SHALL THE COPYRIGHT
      OWNER OR CONTRIBUTORS BE LIABLE FOR ANY DIRECT, INDIRECT, INCIDENTAL,
      SPECIAL, EXEMPLARY, OR CONSEQUENTIAL DAMAGES (INCLUDING, BUT NOT LIMITED
      TO, PROCUREMENT OF SUBSTITUTE GOODS OR SERVICES; LOSS OF USE, DATA, OR
      PROFITS; OR BUSINESS INTERRUPTION) HOWEVER CAUSED AND ON ANY THEORY OF
      LIABILITY, WHETHER IN CONTRACT, STRICT LIABILITY, OR TORT (INCLUDING
      NEGLIGENCE OR OTHERWISE) ARISING IN ANY WAY OUT OF THE USE OF THIS
      SOFTWARE, EVEN IF ADVISED OF THE POSSIBILITY OF SUCH DAMAGE.
     scipy/spatial/qhull/COPYING.txt
     Qhull, Copyright (c) 1993-2015
     C.B. Barber
     Arlington, MA
     and
     The National Science and Technology Research Center for Computation and
      Visualization of Geometric Structures
      (The Geometry Center)
     University of Minnesota
     email: qhull@qhull.org




     This software includes Qhull from C.B. Barber and The Geometry Center.
     Qhull is copyrighted as noted above. Qhull is free software and may
      be obtained via http from www.qhull.org. It may be freely copied,
      modified, and redistributed under the following conditions:
     1. All copyright notices must remain intact in all files.
     2. A copy of this text file must be distributed along with any copies
      of Qhull that you redistribute; this includes copies that you have
      modified, or copies of programs or other software products that include
      Qhull.
     3. If you modify Qhull, you must include a notice giving the name of the
      person performing the modification, the date of modification, and the
      reason for such modification.
     4. When distributing modified versions of Qhull, or other software
      products that include Qhull, you must provide notice that the original
      source code may be obtained as noted above.
     5. There is no warranty or other guarantee of fitness for Qhull, it
      is provided solely "as is". Bug reports or fixes may be sent to
      qhull_bug@qhull.org; the authors may or may not act on them as they
      desire.
\end{lstlisting}

\subsection{Shroud-1.0}
\subsubsection*{Shroud-1.0}

以下为 Shroud-1.0 的许可声明原文。

\begin{lstlisting}[basicstyle=\ttfamily\footnotesize,breaklines=true]
     Project homepage: http://www.cpan.org/authors/id/C/CR/CRAIC/shroud-1.0
     Project license:
     #!/usr/bin/perl
     # shroud
     # Copyright 2000 Robert Jones, Craic Computing, All rights reserved.
     # This program is free software; you can redistribute it and/or modify it
     # under the same terms as Perl itself.
     # The software is supplied as is, with absolutely no warranty.
     Perl5 is Copyright (C) 1993-2005, by Larry Wall and others.

     It is free software; you can redistribute it and/or modify it under the
     terms of either:

     a) the GNU General Public License as published by the Free Software
     Foundation; either external linkversion 1, or (at your option) any later
     versionexternal link, or

     b) the "Artistic License".

     For those of you that choose to use the GNU General Public License, my
     interpretation of the GNU General Public License is that no Perl script
     falls under the terms of the GPL unless you explicitly put said script
     under the terms of the GPL yourself.

     Furthermore, any object code linked with perl does not automatically
     fall under the terms of the GPL, provided such object code only adds
     definitions of subroutines and variables, and does not otherwise impair




     the resulting interpreter from executing any standard Perl script.
     I consider linking in C subroutines in this manner to be the moral
     equivalent of defining subroutines in the Perl language itself. You may
     sell such an object file as proprietary provided that you provide or
     offer to provide the Perl source, as specified by the GNU General Public
     License. (This is merely an alternate way of specifying input to the
     program.) You may also sell a binary produced by the dumping of a running
     Perl script that belongs to you, provided that you provide or offer to
     provide the Perl source as specified by the GPL. (The fact that a Perl
     interpreter and your code are in the same binary file is, in this case, a
     form of mere aggregation.)

     This is my interpretation of the GPL. If you still have concerns or
     difficulties understanding my intent, feel free to contact me. Of course,
     the Artistic License spells all this out for your protection, so you may
     prefer to use that.

     -- Larry Wall

     See Artistic-1.0-Perl in the Standard OSS License Text on page 542."See GPL-1.0 in the
     Standard OSS License Text on page 542."
\end{lstlisting}

\subsection{SQLite}
\subsubsection*{SQLite}

以下为 SQLite 的许可声明原文。

\begin{lstlisting}[basicstyle=\ttfamily\footnotesize,breaklines=true]
     Project homepage/download site: https://www.sqlite.org/index.html
     Project licensing notices: https://www.sqlite.org/copyright.html
     SQLite Is Public Domain

     All of the code and documentation in SQLite has been dedicated to the
     public domain by the authors. All code authors, and representatives
     of the companies they work for, have signed affidavits dedicating
     their contributions to the public domain and originals of those signed
     affidavits are stored in a firesafe at the main offices of Hwaci. Anyone
     is free to copy, modify, publish, use, compile, sell, or distribute the
     original SQLite code, either in source code form or as a compiled binary,
     for any purpose, commercial or non-commercial, and by any means.

     The previous paragraph applies to the deliverable code and documentation
     in SQLite - those parts of the SQLite library that you actually bundle
     and ship with a larger application. Some scripts used as part of the
     build process (for example the "configure" scripts generated by autoconf)
     might fall under other open-source licenses. Nothing from these build
     scripts ever reaches the final deliverable SQLite library, however, and
     so the licenses associated with those scripts should not be a factor in
     assessing your rights to copy and use the SQLite library.




     All of the deliverable code in SQLite has been written from scratch. No
     code has been taken from other projects or from the open internet. Every
     line of code can be traced back to its original author, and all of those
     authors have public domain dedications on file. So the SQLite code base
     is clean and is uncontaminated with licensed code from other projects.

     Open-Source, not Open-Contribution

     SQLite is open-source, meaning that you can make as many copies of it as
     you want and do whatever you want with those copies, without limitation.
     But SQLite is not open-contribution. In order to keep SQLite in the
     public domain and ensure that the code does not become contaminated with
     proprietary or licensed content, the project does not accept patches from
     unknown persons.

     All of the code in SQLite is original, having been written specifically
     for use by SQLite. No code has been copied from unknown sources on the
     internet.

     Warranty of Title

     SQLite is in the public domain and does not require a license. Even
     so, some organizations want legal proof of their right to use SQLite.
     Circumstances where this occurs include the following:

     Your company desires indemnity against claims of copyright infringement.

     You are using SQLite in a jurisdiction that does not recognize the public
     domain.

     You are using SQLite in a jurisdiction that does not recognize the right
     of an author to dedicate their work to the public domain.

     You want to hold a tangible legal document as evidence that you have the
     legal right to use and distribute SQLite.

     Your legal department tells you that you have to purchase a license.

     If any of the above circumstances apply to you, Hwaci, the company that
     employs all the developers of SQLite, will sell you a Warranty of Title
     for SQLite. A Warranty of Title is a legal document that asserts that
     the claimed authors of SQLite are the true authors, and that the authors
     have the legal right to dedicate the SQLite to the public domain, and
     that Hwaci will vigorously defend against challenges to those claims.
     All proceeds from the sale of SQLite Warranties of Title are used to fund
     continuing improvement and support of SQLite.

     Contributed Code




     In order to keep SQLite completely free and unencumbered by copyright,
     the project does not accept patches. If you would like to make a
     suggested change, and include a patch as a proof-of-concept, that would
     be great. However please do not be offended if we rewrite your patch from
     scratch.
\end{lstlisting}

\subsection{TclLib}
\subsubsection*{TclLib}

以下为 TclLib 的许可声明原文。

\begin{lstlisting}[basicstyle=\ttfamily\footnotesize,breaklines=true]
     This software is copyrighted by Ajuba Solutions and other parties.The following terms
     apply to all files associated with the software unless explicitly disclaimed in individual
     files.The authors hereby grant permission to use, copy, modify, distribute, and license this
     software and its documentation for any purpose, provided that existing copyright notices
     are retained in all copies and that this notice is included verbatim in any distributions.
     No written agreement, license, or royalty fee is required for any of the authorized uses.
     Modifications to this software may be copyrighted by their authors and need not follow the
     licensing terms described here, provided that the new terms are clearly indicated on the
     first page of each file where they apply.
     IN NO EVENT SHALL THE AUTHORS OR DISTRIBUTORS BE LIABLE TO ANY
     PARTY FOR DIRECT, INDIRECT, SPECIAL, INCIDENTAL, OR CONSEQUENTIAL
     DAMAGES ARISING OUT OF THE USE OF THIS SOFTWARE, ITS DOCUMENTATION,
     OR ANY DERIVATIVES THEREOF, EVEN IF THE AUTHORS HAVE BEEN ADVISED
     OF THE POSSIBILITY OF SUCH DAMAGE.THE AUTHORS AND DISTRIBUTORS
     SPECIFICALLY DISCLAIM ANY WARRANTIES, INCLUDING, BUT NOT LIMITED TO,
     THE IMPLIED WARRANTIES OF MERCHANTABILITY, FITNESS FOR A PARTICULAR
     PURPOSE, AND NON-INFRINGEMENT. THIS SOFTWARE IS PROVIDED ON AN
     "AS IS" BASIS, AND THE AUTHORS AND DISTRIBUTORS HAVE NO OBLIGATION
     TO PROVIDE MAINTENANCE, SUPPORT, UPDATES, ENHANCEMENTS, OR
     MODIFICATIONS.GOVERNMENT USE: If you are acquiring this software on behalf of
     the U.S. government, the Government shall have only "Restricted Rights" in the software
     and related documentation as defined in the Federal Acquisition Regulations (FARs) in
     Clause 52.227.19 (c) (2). If you are acquiring the software on behalf of the Department
     of Defense, the software shall be classified as "Commercial Computer Software" and the
     Government shall have only "Restricted Rights" as defined in Clause 252.227-7013 (c)
     (1) of DFARs. Notwithstanding the foregoing, the authors grant the U.S. Government and
     others acting in its behalf permission to use and distribute the software in accordance with
     the terms specified in this license.
\end{lstlisting}

\subsection{Tcl/Tk}
\subsubsection*{Tcl/Tk}

以下为 Tcl/Tk 的许可声明原文。

\begin{lstlisting}[basicstyle=\ttfamily\footnotesize,breaklines=true]
     The following terms apply to the all versions of the core Tcl/Tk releases, the Tcl/Tk browser
     plug-in version 2.0, Tcllib, and TclBlend and Jacl version 1.0. Please note that the TclPro
     tools are under a different license agreement. This agreement is part of the standard Tcl/
     Tk distribution as the file named "license.terms".




     Tcl/Tk License Terms
     This software is copyrighted by the Regents of the University of California, Sun
     Microsystems, Inc., Scriptics Corporation, Ajuba Solutions, and other parties. The
     following terms apply to all files associated with the software unless explicitly disclaimed in
     individual files.
     The authors hereby grant permission to use, copy, modify, distribute, and license this
     software and its documentation for any purpose, provided that existing copyright notices
     are retained in all copies and that this notice is included verbatim in any distributions.
     No written agreement, license, or royalty fee is required for any of the authorized uses.
     Modifications to this software may be copyrighted by their authors and need not follow the
     licensing terms described here, provided that the new terms are clearly indicated on the
     first page of each file where they apply.
     IN NO EVENT SHALL THE AUTHORS OR DISTRIBUTORS BE LIABLE TO ANY PARTY
     FOR DIRECT, INDIRECT, SPECIAL, INCIDENTAL, OR CONSEQUENTIAL DAMAGES
     ARISING OUT OF THE USE OF THIS SOFTWARE, ITS DOCUMENTATION, OR ANY
     DERIVATIVES THEREOF, EVEN IF THE AUTHORS HAVE BEEN ADVISED OF THE
     POSSIBILITY OF SUCH DAMAGE.
     THE AUTHORS AND DISTRIBUTORS SPECIFICALLY DISCLAIM ANY
     WARRANTIES, INCLUDING, BUT NOT LIMITED TO, THE IMPLIED WARRANTIES
     OF MERCHANTABILITY, FITNESS FOR A PARTICULAR PURPOSE, AND NON-
     INFRINGEMENT. THIS SOFTWARE IS PROVIDED ON AN "AS IS" BASIS, AND
     THE AUTHORS AND DISTRIBUTORS HAVE NO OBLIGATION TO PROVIDE
     MAINTENANCE, SUPPORT, UPDATES, ENHANCEMENTS, OR MODIFICATIONS.
     GOVERNMENT USE: If you are acquiring this software on behalf of the U.S. government,
     the Government shall have only "Restricted Rights" in the software and related
     documentation as defined in the Federal Acquisition Regulations (FARs) in Clause
     52.227.19 (c) (2). If you are acquiring the software on behalf of the Department of
     Defense, the software shall be classified as "Commercial Computer Software" and the
     Government shall have only "Restricted Rights" as defined in Clause 252.227-7013 (c)
     (1) of DFARs. Notwithstanding the foregoing, the authors grant the U.S. Government and
     others acting in its behalf permission to use and distribute the software in accordance with
     the terms specified in this license.
\end{lstlisting}

\subsection{zlib}
\subsubsection*{zlib}

以下为 zlib 的许可声明原文。

\begin{lstlisting}[basicstyle=\ttfamily\footnotesize,breaklines=true]
     License
     /* zlib.h -- interface of the 'zlib' general purpose compression library version 1.2.8, April
     28th, 2013
     Copyright (C) 1995-2013 Jean-loup Gailly and Mark Adler




     This software is provided 'as-is', without any express or implied warranty. In no event will
     the authors be held liable for any damages arising from the use of this software.
     Permission is granted to anyone to use this software for any purpose, including
     commercial applications, and to alter it and redistribute it freely, subject to the following
     restrictions:
      1. The origin of this software must not be misrepresented; you must not claim that you
         wrote the original software. If you use this software in a product, an acknowledgement
         in the product documentation would be appreciated but is not required.
      2. Altered source versions must be plainly marked as such, and must not be
         misrepresented as being the original software.
      3. This notice may not be removed or altered from any source distribution.
     Jean-loup Gailly (jloup@gzip.org)       Mark Adler (madler@alumni.caltech.edu)
     */
\end{lstlisting}

\subsection{Zstandard}
\subsubsection*{Zstandard}

以下为 Zstandard 的许可声明原文。

\begin{lstlisting}[basicstyle=\ttfamily\footnotesize,breaklines=true]
     Project homepage: https://github.com/facebook/zstd
     Project license: https://github.com/facebook/zstd/blob/master/LICENSE
     BSD License

     For Zstandard software

     Copyright (c) 2016-present, Facebook, Inc. All rights reserved.

     Redistribution and use in source and binary forms, with or without
     modification, are permitted provided that the following conditions are
     met:

     * Redistributions of source code must retain the above copyright notice,
     this list of conditions and the following disclaimer.

     * Redistributions in binary form must reproduce the above copyright
     notice, this list of conditions and the following disclaimer in the
     documentation and/or other materials provided with the distribution.

     * Neither the name Facebook nor the names of its contributors may be
     used to endorse or promote products derived from this software without
     specific prior written permission.

     THIS SOFTWARE IS PROVIDED BY THE COPYRIGHT HOLDERS AND CONTRIBUTORS
     "AS IS" AND ANY EXPRESS OR IMPLIED WARRANTIES, INCLUDING, BUT NOT
     LIMITED TO, THE IMPLIED WARRANTIES OF MERCHANTABILITY AND FITNESS FOR
\end{lstlisting}

% ============================================================
\section{标准 OSS 许可文本}
\subsection*{Standard OSS License Text}

本节转载若干标准开源许可的完整法律文本（英文原文）。

\subsection{Apache-2.0}
\subsubsection*{Apache-2.0}

以下为 Apache-2.0 许可原文。

\begin{lstlisting}[basicstyle=\ttfamily\footnotesize,breaklines=true]
     Project homepage/download site: http://google.github.io/flatbuffers/; https://github.com/
     google/flatbuffers
     Project licensing notices: https://github.com/google/flatbuffers/blob/master/LICENSE.txt
     LICENSE.txt
     Apache LicenseVersion 2.0, January 2004 http://www.apache.org/licenses/

     TERMS AND CONDITIONS FOR USE, REPRODUCTION, AND DISTRIBUTION

     1. Definitions.

     "License" shall mean the terms and conditions for use, reproduction,and
     distribution as defined by Sections 1 through 9 of this document.

     "Licensor" shall mean the copyright owner or entity authorized by the
     copyright owner that is granting the License.

     "Legal Entity" shall mean the union of the acting entity and all other
     entities that control, are controlled by, or are under common control
     with that entity. For the purposes of this definition, "control" means
     (i) the power, direct or indirect, to cause the direction or management
     of such entity, whether by contract or otherwise, or (ii) ownership
     of fifty percent (50%) or more of the outstanding shares, or (iii)
     beneficial ownership of such entity.

     "You" (or "Your") shall mean an individual or Legal Entity exercising
     permissions granted by this License.

     "Source" form shall mean the preferred form for making modifications,
     including but not limited to software source code, documentation source,
     and configuration files.




     "Object" form shall mean any form resulting from mechanical
     transformation or translation of a Source form, including but not limited
     to compiled object code, generated documentation, and conversions to
     other media types.

     "Work" shall mean the work of authorship, whether in Source or Object
     form, made available under the License, as indicated by a copyright
     notice that is included in or attached to the work (an example is
     provided in the Appendix below).

     "Derivative Works" shall mean any work, whether in Source or Object form,
     that is based on (or derived from) the Work and for which the editorial
     revisions, annotations, elaborations, or other modifications represent,
     as a whole, an original work of authorship. For the purposes of this
     License, Derivative Works shall not include works that remain separable
     from, or merely link (or bind by name) to the interfaces of, the Work and
     Derivative Works thereof.

     "Contribution" shall mean any work of authorship, including the original
     version of the Work and any modifications or additions to that Work or
     Derivative Works thereof, that is intentionally submitted to Licensor
     for inclusion in the Work by the copyright owner or by an individual
     or Legal Entity authorized to submit on behalf of the copyright owner.
     For the purposes of this definition, "submitted" means any form of
     electronic, verbal, or written communication sent to the Licensor or its
     representatives, including but not limited to communication on electronic
     mailing lists, source code control systems, and issue tracking systems
     that are managed by, or on behalf of, the Licensor for the purpose of
     discussing and improving the Work, but excluding communication that is
     conspicuously marked or otherwise designated in writing by the copyright
     owner as "Not a Contribution."

     "Contributor" shall mean Licensor and any individual or Legal Entity
     on behalf of whom a Contribution has been received by Licensor and
     subsequently incorporated within the Work.

     2. Grant of Copyright License. Subject to the terms and conditions
     of this License, each Contributor hereby grants to You a perpetual,
     worldwide, non-exclusive, no-charge, royalty-free, irrevocable copyright
     license to reproduce, prepare Derivative Works of, publicly display,
     publicly perform, sublicense, and distribute the Work and such Derivative
     Works in Source or Object form.

     3. Grant of Patent License. Subject to the terms and conditions of this
     License, each Contributor hereby grants to You a perpetual, worldwide,
     non-exclusive, no-charge, royalty-free, irrevocable (except as stated
     in this section) patent license to make, have made, use, offer to sell,




     sell, import, and otherwise transfer the Work, where such license applies
     only to those patent claims licensable by such Contributor that are
     necessarily infringed by their Contribution(s) alone or by combination
     of their Contribution(s) with the Work to which such Contribution(s)
     was submitted. If You institute patent litigation against any entity
     (including a cross-claim or counterclaim in a lawsuit) alleging that the
     Work or a Contribution incorporated within the Work constitutes direct
     or contributory patent infringement, then any patent licenses granted to
     You under this License for that Work shall terminate as of the date such
     litigation is filed.

     4. Redistribution. You may reproduce and distribute copies of the Work or
     Derivative Works thereof in any medium, with or without modifications,
     and in Source or Object form, provided that You meet the following
     conditions:

     (a) You must give any other recipients of the Work or Derivative Works a
     copy of this License; and

     (b) You must cause any modified files to carry prominent notices stating
     that You changed the files; and

     (c) You must retain, in the Source form of any Derivative Works that You
     distribute, all copyright, patent, trademark, and attribution notices
     from the Source form of the Work, excluding those notices that do not
     pertain to any part of the Derivative Works; and

     (d) If the Work includes a "NOTICE" text file as part of its
     distribution, then any Derivative Works that You distribute must include
     a readable copy of the attribution notices contained within such NOTICE
     file, excluding those notices that do not pertain to any part of the
     Derivative Works, in at least one of the following places: within a
     NOTICE text file distributed as part of the Derivative Works; within
     the Source form or documentation, if provided along with the Derivative
     Works; or, within a display generated by the Derivative Works, if and
     wherever such third-party notices normally appear. The contents of the
     NOTICE file are for informational purposes only and do not modify the
     License. You may add Your own attribution notices within Derivative Works
     that You distribute, alongside or as an addendum to the NOTICE text from
     the Work, provided that such additional attribution notices cannot be
     construed as modifying the License.

     You may add Your own copyright statement to Your modifications and
     may provide additional or different license terms and conditions for
     use, reproduction, or distribution of Your modifications, or for any
     such Derivative Works as a whole, provided Your use, reproduction, and




     distribution of the Work otherwise complies with the conditions stated in
     this License.

     5. Submission of Contributions. Unless You explicitly state otherwise,
     any Contribution intentionally submitted for inclusion in the Work by
     You to the Licensor shall be under the terms and conditions of this
     License, without any additional terms or conditions. Notwithstanding the
     above, nothing herein shall supersede or modify the terms of any separate
     license agreement you may have executed with Licensor regarding such
     Contributions.

     6. Trademarks. This License does not grant permission to use the trade
     names, trademarks, service marks, or product names of the Licensor,
     except as required for reasonable and customary use in describing the
     origin of the Work and reproducing the content of the NOTICE file.

     7. Disclaimer of Warranty. Unless required by applicable law or agreed
     to in writing, Licensor provides the Work (and each Contributor provides
     its Contributions) on an "AS IS" BASIS, WITHOUT WARRANTIES OR CONDITIONS
     OF ANY KIND, either express or implied, including, without limitation,
     any warranties or conditions of TITLE, NON-INFRINGEMENT, MERCHANTABILITY,
     or FITNESS FOR A PARTICULAR PURPOSE. You are solely responsible for
     determining the appropriateness of using or redistributing the Work and
     assume any risks associated with Your exercise of permissions under this
     License.

     8. Limitation of Liability. In no event and under no legal theory,
     whether in tort (including negligence), contract, or otherwise, unless
     required by applicable law (such as deliberate and grossly negligent
     acts) or agreed to in writing, shall any Contributor be liable to You
     for damages, including any direct, indirect, special, incidental, or
     consequential damages of any character arising as a result of this
     License or out of the use or inability to use the Work (including but not
     limited to damages for loss of goodwill, work stoppage, computer failure
     or malfunction, or any and all other commercial damages or losses), even
     if such Contributor has been advised of the possibility of such damages.

     9. Accepting Warranty or Additional Liability. While redistributing the
     Work or Derivative Works thereof, You may choose to offer, and charge a
     fee for, acceptance of support, warranty, indemnity, or other liability
     obligations and/or rights consistent with this License. However, in
     accepting such obligations, You may act only on Your own behalf and
     on Your sole responsibility, not on behalf of any other Contributor,
     and only if You agree to indemnify, defend, and hold each Contributor
     harmless for any liability incurred by, or claims asserted against, such




     Contributor by reason of your accepting any such warranty or additional
     liability.

     END OF TERMS AND CONDITIONS\

     APPENDIX: How to apply the Apache License to your work.

     To apply the Apache License to your work, attach the following
     boilerplate notice, with the fields enclosed by brackets "[]" replaced
     with your own identifying information. (Don't include the brackets!) The
     text should be enclosed in the appropriate comment syntax for the file
     format. We also recommend that a file or class name and description of
     purpose be included on the same "printed page" as the copyright notice
     for easier identification within third-party archives.

     Copyright 2014 Google Inc.

     Licensed under the Apache License, Version 2.0 (the "License"); you may
     not use this file except in compliance with the License. You may obtain a
     copy of the License at

     http://www.apache.org/licenses/LICENSE-2.0

     Unless required by applicable law or agreed to in writing, software
     distributed under the License is distributed on an "AS IS" BASIS, WITHOUT
     WARRANTIES OR CONDITIONS OF ANY KIND, either express or implied. See the
     License for the specific language governing permissions and limitations
     under the License.
\end{lstlisting}

\subsection{Artistic-1.0-Perl}
\subsubsection*{Artistic-1.0-Perl}

以下为 Artistic-1.0-Perl 许可原文。

\begin{lstlisting}[basicstyle=\ttfamily\footnotesize,breaklines=true]
     The "Artistic License"

     Preamble

     The intent of this document is to state the conditions under which a
     Package may be copied, such that the Copyright Holder maintains some
     semblance of artistic control over the development of the package, while
     giving the users of the package the right to use and distribute the
     Package in a more-or-less customary fashion, plus the right to make
     reasonable modifications.

     Definitions:

     "Package" refers to the collection of files distributed by the Copyright
     Holder, and derivatives of that collection of files created through
     textual modification.




     "Standard Version" refers to such a Package if it has not been modified,
     or has been modified in accordance with the wishes of the Copyright
     Holder as specified below.

     "Copyright Holder" is whoever is named in the copyright or copyrights for
     the package.

     "You" is you, if you're thinking about copying or distributing this
     Package.

     "Reasonable copying fee" is whatever you can justify on the basis of
     media cost, duplication charges, time of people involved, and so on. (You
     will not be required to justify it to the Copyright Holder, but only to
     the computing community at large as a market that must bear the fee.)

     "Freely Available" means that no fee is charged for the item itself,
     though there may be fees involved in handling the item. It also means
     that recipients of the item may redistribute it under the same conditions
     they received it.

      • You may make and give away verbatim copies of the source form of the
         Standard Version of this Package without restriction, provided that
         you duplicate all of the original copyright notices and associated
         disclaimers.

      • You may apply bug fixes, portability fixes and other modifications
         derived from the Public Domain or from the Copyright Holder. A Package
         modified in such a way shall still be considered the Standard Version.

      • You may otherwise modify your copy of this Package in any way,
         provided that you insert a prominent notice in each changed file
         stating how and when you changed that file, and provided that you do
         at least ONE of the following:

          ◦ place your modifications in the Public Domain or otherwise make them
             Freely Available, such as by posting said modifications to Usenet
             or an equivalent medium, or placing the modifications on a major
             archive site such as uunet.uu.net, or by allowing the Copyright
             Holder to include your modifications in the Standard Version of the
             Package.

          ◦ use the modified Package only within your corporation or
             organization.

          ◦ rename any non-standard executables so the names do not conflict
             with standard executables, which must also be provided, and provide




             a separate manual page for each non-standard executable that clearly
             documents how it differs from the Standard Version.

          ◦ make other distribution arrangements with the Copyright Holder.
      • You may distribute the programs of this Package in object code or
         executable form, provided that you do at least ONE of the following:

          ◦ distribute a Standard Version of the executables and library files,
             together with instructions (in the manual page or equivalent) on
             where to get the Standard Version.

          ◦ accompany the distribution with the machine-readable source of the
             Package with your modifications.

          ◦ give non-standard executables non-standard names, and clearly
             document the differences in manual pages (or equivalent), together
             with instructions on where to get the Standard Version.

          ◦ make other distribution arrangements with the Copyright Holder.
      • You may charge a reasonable copying fee for any distribution of
         this Package. You may charge any fee you choose for support of this
         Package. You may not charge a fee for this Package itself. However,
         you may distribute this Package in aggregate with other (possibly
         commercial) programs as part of a larger (possibly commercial)
         software distribution provided that you do not advertise this Package
         as a product of your own. You may embed this Package's interpreter
         within an executable of yours (by linking); this shall be construed
         as a mere form of aggregation, provided that the complete Standard
         Version of the interpreter is so embedded.

      • The scripts and library files supplied as input to or produced as
         output from the programs of this Package do not automatically fall
         under the copyright of this Package, but belong to whoever generated
         them, and may be sold commercially, and may be aggregated with this
         Package. If such scripts or library files are aggregated with this
         Package via the so-called "undump" or "unexec" methods of producing
         a binary executable image, then distribution of such an image shall
         neither be construed as a distribution of this Package nor shall it
         fall under the restrictions of Paragraphs 3 and 4, provided that you
         do not represent such an executable image as a Standard Version of
         this Package.

      • subroutines (or comparably compiled subroutines in other languages)
         supplied by you and linked into this Package in order to emulate
         subroutines and variables of the language defined by this Package
         shall not be considered part of this Package, but are the equivalent




         of input as in Paragraph 6, provided these subroutines do not change
         the language in any way that would cause it to fail the regression
         tests for the language.

      • Aggregation of this Package with a commercial distribution is always
         permitted provided that the use of this Package is embedded; that
         is, when no overt attempt is made to make this Package's interfaces
         visible to the end user of the commercial distribution. Such use shall
         not be construed as a distribution of this Package.

      • The name of the Copyright Holder may not be used to endorse or promote
         products derived from this software without specific prior written
         permission.

      • THIS PACKAGE IS PROVIDED "AS IS" AND WITHOUT ANY EXPRESS OR IMPLIED
         WARRANTIES, INCLUDING, WITHOUT LIMITATION, THE IMPLIED WARRANTIES OF
         MERCHANTIBILITY AND FITNESS FOR A PARTICULAR PURPOSE.

     The End
\end{lstlisting}

\subsection{GPL-1.0}
\subsubsection*{GPL-1.0}

以下为 GPL-1.0 许可原文。

\begin{lstlisting}[basicstyle=\ttfamily\footnotesize,breaklines=true]
     GNU GENERAL PUBLIC LICENSE Version 1, February 1989

     Copyright (C) 1989 Free Software Foundation, Inc. 59 Temple Place,
     Suite 330, Boston, MA 02111-1307, USA Everyone is permitted to copy and
     distribute verbatim copies of this license document, but changing it is
     not allowed.

     Preamble

     The license agreements of most software companies try to keep users at
     the mercy of those companies. By contrast, our General Public License is
     intended to guarantee your freedom to share and change free software--
     to make sure the software is free for all its users. The General Public
     License applies to the Free Software Foundation's software and to any
     other program whose authors commit to using it. You can use it for your
     programs, too.

     When we speak of free software, we are referring to freedom, not price.
     Specifically, the General Public License is designed to make sure that
     you have the freedom to give away or sell copies of free software, that
     you receive source code or can get it if you want it, that you can change
     the software or use pieces of it in new free programs; and that you know
     you can do these things.




     To protect your rights, we need to make restrictions that forbid anyone
     to deny you these rights or to ask you to surrender the rights. These
     restrictions translate to certain responsibilities for you if you
     distribute copies of the software, or if you modify it.

     For example, if you distribute copies of a such a program, whether gratis
     or for a fee, you must give the recipients all the rights that you have.
     You must make sure that they, too, receive or can get the source code.
     And you must tell them their rights.

     We protect your rights with two steps: (1) copyright the software, and
     (2) offer you this license which gives you legal permission to copy,
     distribute and/or modify the software.

     Also, for each author's protection and ours, we want to make certain that
     everyone understands that there is no warranty for this free software.
     If the software is modified by someone else and passed on, we want its
     recipients to know that what they have is not the original, so that any
     problems introduced by others will not reflect on the original authors'
     reputations.

     The precise terms and conditions for copying, distribution and
     modification follow. GNU GENERAL PUBLIC LICENSE TERMS AND CONDITIONS FOR
     COPYING, DISTRIBUTION AND MODIFICATION

     0. This License Agreement applies to any program or other work which
     contains a notice placed by the copyright holder saying it may be
     distributed under the terms of this General Public License. The
     "Program", below, refers to any such program or work, and a "work based
     on the Program" means either the Program or any work containing the
     Program or a portion of it, either verbatim or with modifications. Each
     licensee is addressed as "you".

      • You may copy and distribute verbatim copies of the Program's source
         code as you receive it, in any medium, provided that you conspicuously
         and appropriately publish on each copy an appropriate copyright notice
         and disclaimer of warranty; keep intact all the notices that refer to
         this General Public License and to the absence of any warranty; and
         give any other recipients of the Program a copy of this General Public
         License along with the Program. You may charge a fee for the physical
         act of transferring a copy.




      • You may modify your copy or copies of the Program or any portion of
         it, and copy and distribute such modifications under the terms of
         Paragraph 1 above, provided that you also do the following:

          ◦ cause the modified files to carry prominent notices stating that you
             changed the files and the date of any change; and

          ◦ cause the whole of any work that you distribute or publish, that in
             whole or in part contains the Program or any part thereof, either
             with or without modifications, to be licensed at no charge to all
             third parties under the terms of this General Public License (except
             that you may choose to grant warranty protection to some or all
             third parties, at your option).

          ◦ If the modified program normally reads commands interactively when
             run, you must cause it, when started running for such interactive
             use in the simplest and most usual way, to print or display an
             announcement including an appropriate copyright notice and a notice
             that there is no warranty (or else, saying that you provide a
             warranty) and that users may redistribute the program under these
             conditions, and telling the user how to view a copy of this General
             Public License.

          ◦ You may charge a fee for the physical act of transferring a copy,
             and you may at your option offer warranty protection in exchange for
             a fee.

             Mere aggregation of another independent work with the Program (or
             its derivative) on a volume of a storage or distribution medium does
             not bring the other work under the scope of these terms. 3. You may
             copy and distribute the Program (or a portion or derivative of it,
             under Paragraph 2) in object code or executable form under the terms
             of Paragraphs 1 and 2 above provided that you also do one of the
             following:

          ◦ accompany it with the complete corresponding machine-readable source
             code, which must be distributed under the terms of Paragraphs 1 and
             2 above; or,

          ◦ accompany it with a written offer, valid for at least three years,
             to give any third party free (except for a nominal charge for
             the cost of distribution) a complete machine-readable copy of the
             corresponding source code, to be distributed under the terms of
             Paragraphs 1 and 2 above; or,

          ◦ accompany it with the information you received as to where the
             corresponding source code may be obtained. (This alternative is




             allowed only for noncommercial distribution and only if you received
             the program in object code or executable form alone.)

             Source code for a work means the preferred form of the work for
             making modifications to it. For an executable file, complete source
             code means all the source code for all modules it contains; but,
             as a special exception, it need not include source code for modules
             which are standard libraries that accompany the operating system
             on which the executable file runs, or for standard header files or
             definitions files that accompany that operating system.

     4. You may not copy, modify, sublicense, distribute or transfer the
     Program except as expressly provided under this General Public License.
     Any attempt otherwise to copy, modify, sublicense, distribute or transfer
     the Program is void, and will automatically terminate your rights to
     use the Program under this License. However, parties who have received
     copies, or rights to use copies, from you under this General Public
     License will not have their licenses terminated so long as such parties
     remain in full compliance.

     5. By copying, distributing or modifying the Program (or any work based on
     the Program) you indicate your acceptance of this license to do so, and
     all its terms and conditions.

     6. Each time you redistribute the Program (or any work based on the
     Program), the recipient automatically receives a license from the
     original licensor to copy, distribute or modify the Program subject to
     these terms and conditions. You may not impose any further restrictions
     on the recipients' exercise of the rights granted herein. 7. The Free
     Software Foundation may publish revised and/or new versions of the
     General Public License from time to time. Such new versions will be
     similar in spirit to the present version, but may differ in detail to
     address new problems or concerns.

     Each version is given a distinguishing version number. If the Program
     specifies a version number of the license which applies to it and "any
     later version", you have the option of following the terms and conditions
     either of that version or of any later version published by the Free
     Software Foundation. If the Program does not specify a version number
     of the license, you may choose any version ever published by the Free
     Software Foundation.

     8. If you wish to incorporate parts of the Program into other free
     programs whose distribution conditions are different, write to the author
     to ask for permission. For software which is copyrighted by the Free
     Software Foundation, write to the Free Software Foundation; we sometimes
     make exceptions for this. Our decision will be guided by the two goals of




     preserving the free status of all derivatives of our free software and of
     promoting the sharing and reuse of software generally.

     NO WARRANTY

     9. BECAUSE THE PROGRAM IS LICENSED FREE OF CHARGE, THERE IS NO WARRANTY
     FOR THE PROGRAM, TO THE EXTENT PERMITTED BY APPLICABLE LAW. EXCEPT
     WHEN OTHERWISE STATED IN WRITING THE COPYRIGHT HOLDERS AND/OR OTHER
     PARTIES PROVIDE THE PROGRAM "AS IS" WITHOUT WARRANTY OF ANY KIND,
     EITHER EXPRESSED OR IMPLIED, INCLUDING, BUT NOT LIMITED TO, THE IMPLIED
     WARRANTIES OF MERCHANTABILITY AND FITNESS FOR A PARTICULAR PURPOSE. THE
     ENTIRE RISK AS TO THE QUALITY AND PERFORMANCE OF THE PROGRAM IS WITH YOU.
     SHOULD THE PROGRAM PROVE DEFECTIVE, YOU ASSUME THE COST OF ALL NECESSARY
     SERVICING, REPAIR OR CORRECTION.

     10. IN NO EVENT UNLESS REQUIRED BY APPLICABLE LAW OR AGREED TO IN
     WRITING WILL ANY COPYRIGHT HOLDER, OR ANY OTHER PARTY WHO MAY MODIFY AND/
     OR REDISTRIBUTE THE PROGRAM AS PERMITTED ABOVE, BE LIABLE TO YOU FOR
     DAMAGES, INCLUDING ANY GENERAL, SPECIAL, INCIDENTAL OR CONSEQUENTIAL
     DAMAGES ARISING OUT OF THE USE OR INABILITY TO USE THE PROGRAM (INCLUDING
     BUT NOT LIMITED TO LOSS OF DATA OR DATA BEING RENDERED INACCURATE OR
     LOSSES SUSTAINED BY YOU OR THIRD PARTIES OR A FAILURE OF THE PROGRAM TO
     OPERATE WITH ANY OTHER PROGRAMS), EVEN IF SUCH HOLDER OR OTHER PARTY HAS
     BEEN ADVISED OF THE POSSIBILITY OF SUCH DAMAGES.

     END OF TERMS AND CONDITIONS Appendix: How to Apply These Terms to Your
     New Programs

     If you develop a new program, and you want it to be of the greatest
     possible use to humanity, the best way to achieve this is to make it free
     software which everyone can redistribute and change under these terms.

     To do so, attach the following notices to the program. It is safest
     to attach them to the start of each source file to most effectively
     convey the exclusion of warranty; and each file should have at least the
     "copyright" line and a pointer to where the full notice is found.

     Copyright (C) 19yy

     This program is free software; you can redistribute it and/or modify it
     under the terms of the GNU General Public License as published by the
     Free Software Foundation; either version 1, or (at your option) any later
     version.

     This program is distributed in the hope that it will be useful,
     but WITHOUT ANY WARRANTY; without even the implied warranty of
     MERCHANTABILITY or FITNESS FOR A PARTICULAR PURPOSE. See the GNU General
     Public License for more details.




     You should have received a copy of the GNU General Public License along
     with this program; if not, write to the Free Software Foundation, Inc.,
     59 Temple Place, Suite 330, Boston, MA 02111-1307, USA.

     Also add information on how to contact you by electronic and paper mail.

     If the program is interactive, make it output a short notice like this
     when it starts in an interactive mode:

     Gnomovision version 69, Copyright (C) 19xx name of author Gnomovision
     comes with ABSOLUTELY NO WARRANTY; for details type `show w'. This is
     free software, and you are welcome to redistribute it under certain
     conditions; type `show c' for details.

     The hypothetical commands `show w' and `show c' should show the
     appropriate parts of the General Public License. Of course, the commands
     you use may be called something other than `show w' and `show c'; they
     could even be mouse-clicks or menu items--whatever suits your program.

     You should also get your employer (if you work as a programmer) or your
     school, if any, to sign a "copyright disclaimer" for the program, if
     necessary. Here a sample; alter the names:

     Yoyodyne, Inc., hereby disclaims all copyright interest in the
     program `Gnomovision' (a program to direct compilers to make passes at
     assemblers) written by James Hacker.

     , 1 April 1989 Ty Coon, President of Vice

     That's all there is to it!
\end{lstlisting}

\subsection{GPL-3.0}
\subsubsection*{GPL-3.0}

以下为 GPL-3.0 许可原文。

\begin{lstlisting}[basicstyle=\ttfamily\footnotesize,breaklines=true]
     GNU GENERAL PUBLIC LICENSE Version 3, 29 June 2007

     Copyright (C) 2007 Free Software Foundation, Inc. <http://fsf.org/>
     Everyone is permitted to copy and distribute verbatim copies of this
     license document, but changing it is not allowed.

     Preamble

     The GNU General Public License is a free, copyleft license for software
     and other kinds of works.

     The licenses for most software and other practical works are designed
     to take away your freedom to share and change the works. By contrast,
     the GNU General Public License is intended to guarantee your freedom to
     share and change all versions of a program--to make sure it remains free
     software for all its users. We, the Free Software Foundation, use the




     GNU General Public License for most of our software; it applies also to
     any other work released this way by its authors. You can apply it to your
     programs, too.

     When we speak of free software, we are referring to freedom, not price.
     Our General Public Licenses are designed to make sure that you have the
     freedom to distribute copies of free software (and charge for them if you
     wish), that you receive source code or can get it if you want it, that
     you can change the software or use pieces of it in new free programs, and
     that you know you can do these things.

     To protect your rights, we need to prevent others from denying you these
     rights or asking you to surrender the rights. Therefore, you have certain
     responsibilities if you distribute copies of the software, or if you
     modify it: responsibilities to respect the freedom of others.

     For example, if you distribute copies of such a program, whether gratis
     or for a fee, you must pass on to the recipients the same freedoms that
     you received. You must make sure that they, too, receive or can get
     the source code. And you must show them these terms so they know their
     rights.

     Developers that use the GNU GPL protect your rights with two steps: (1)
     assert copyright on the software, and (2) offer you this License giving
     you legal permission to copy, distribute and/or modify it.

     For the developers' and authors' protection, the GPL clearly explains
     that there is no warranty for this free software. For both users' and
     authors' sake, the GPL requires that modified versions be marked as
     changed, so that their problems will not be attributed erroneously to
     authors of previous versions.

     Some devices are designed to deny users access to install or run modified
     versions of the software inside them, although the manufacturer can do
     so. This is fundamentally incompatible with the aim of protecting users'
     freedom to change the software. The systematic pattern of such abuse
     occurs in the area of products for individuals to use, which is precisely
     where it is most unacceptable. Therefore, we have designed this version
     of the GPL to prohibit the practice for those products. If such problems
     arise substantially in other domains, we stand ready to extend this
     provision to those domains in future versions of the GPL, as needed to
     protect the freedom of users.

     Finally, every program is threatened constantly by software patents.
     States should not allow patents to restrict development and use of
     software on general-purpose computers, but in those that do, we wish to
     avoid the special danger that patents applied to a free program could




     make it effectively proprietary. To prevent this, the GPL assures that
     patents cannot be used to render the program non-free.

     The precise terms and conditions for copying, distribution and
     modification follow.

     TERMS AND CONDITIONS

     0. Definitions.

     "This License" refers to version 3 of the GNU General Public License.

     "Copyright" also means copyright-like laws that apply to other kinds of
     works, such as semiconductor masks.

     "The Program" refers to any copyrightable work licensed under
     this License. Each licensee is addressed as "you". "Licensees" and
     "recipients" may be individuals or organizations.

     To "modify" a work means to copy from or adapt all or part of the work
     in a fashion requiring copyright permission, other than the making of
     an exact copy. The resulting work is called a "modified version" of the
     earlier work or a work "based on" the earlier work.

     A "covered work" means either the unmodified Program or a work based on
     the Program.

     To "propagate" a work means to do anything with it that, without
     permission, would make you directly or secondarily liable for
     infringement under applicable copyright law, except executing it on a
     computer or modifying a private copy. Propagation includes copying,
     distribution (with or without modification), making available to the
     public, and in some countries other activities as well.

     To "convey" a work means any kind of propagation that enables other
     parties to make or receive copies. Mere interaction with a user through a
     computer network, with no transfer of a copy, is not conveying.

     An interactive user interface displays "Appropriate Legal Notices" to the
     extent that it includes a convenient and prominently visible feature that
     (1) displays an appropriate copyright notice, and (2) tells the user that
     there is no warranty for the work (except to the extent that warranties
     are provided), that licensees may convey the work under this License, and
     how to view a copy of this License. If the interface presents a list of
     user commands or options, such as a menu, a prominent item in the list
     meets this criterion.

     1. Source Code.




     The "source code" for a work means the preferred form of the work for
     making modifications to it. "Object code" means any non-source form of a
     work.

     A "Standard Interface" means an interface that either is an official
     standard defined by a recognized standards body, or, in the case of
     interfaces specified for a particular programming language, one that is
     widely used among developers working in that language.

     The "System Libraries" of an executable work include anything, other than
     the work as a whole, that (a) is included in the normal form of packaging
     a Major Component, but which is not part of that Major Component, and (b)
     serves only to enable use of the work with that Major Component, or to
     implement a Standard Interface for which an implementation is available
     to the public in source code form. A "Major Component", in this context,
     means a major essential component (kernel, window system, and so on) of
     the specific operating system (if any) on which the executable work runs,
     or a compiler used to produce the work, or an object code interpreter
     used to run it.

     The "Corresponding Source" for a work in object code form means all
     the source code needed to generate, install, and (for an executable
     work) run the object code and to modify the work, including scripts to
     control those activities. However, it does not include the work's System
     Libraries, or general-purpose tools or generally available free programs
     which are used unmodified in performing those activities but which
     are not part of the work. For example, Corresponding Source includes
     interface definition files associated with source files for the work, and
     the source code for shared libraries and dynamically linked subprograms
     that the work is specifically designed to require, such as by intimate
     data communication or control flow between those subprograms and other
     parts of the work.

     The Corresponding Source need not include anything that users can
     regenerate automatically from other parts of the Corresponding Source.

     The Corresponding Source for a work in source code form is that same
     work.

     2. Basic Permissions.

     All rights granted under this License are granted for the term of
     copyright on the Program, and are irrevocable provided the stated
     conditions are met. This License explicitly affirms your unlimited
     permission to run the unmodified Program. The output from running a
     covered work is covered by this License only if the output, given its




     content, constitutes a covered work. This License acknowledges your
     rights of fair use or other equivalent, as provided by copyright law.

     You may make, run and propagate covered works that you do not convey,
     without conditions so long as your license otherwise remains in force.
     You may convey covered works to others for the sole purpose of having
     them make modifications exclusively for you, or provide you with
     facilities for running those works, provided that you comply with the
     terms of this License in conveying all material for which you do not
     control copyright. Those thus making or running the covered works for you
     must do so exclusively on your behalf, under your direction and control,
     on terms that prohibit them from making any copies of your copyrighted
     material outside their relationship with you.

     Conveying under any other circumstances is permitted solely under the
     conditions stated below. Sublicensing is not allowed; section 10 makes it
     unnecessary.

     3. Protecting Users' Legal Rights From Anti-Circumvention Law.

     No covered work shall be deemed part of an effective technological
     measure under any applicable law fulfilling obligations under article 11
     of the WIPO copyright treaty adopted on 20 December 1996, or similar laws
     prohibiting or restricting circumvention of such measures.

     When you convey a covered work, you waive any legal power to forbid
     circumvention of technological measures to the extent such circumvention
     is effected by exercising rights under this License with respect to
     the covered work, and you disclaim any intention to limit operation or
     modification of the work as a means of enforcing, against the work's
     users, your or third parties' legal rights to forbid circumvention of
     technological measures.

     4. Conveying Verbatim Copies.

     You may convey verbatim copies of the Program's source code as
     you receive it, in any medium, provided that you conspicuously and
     appropriately publish on each copy an appropriate copyright notice; keep
     intact all notices stating that this License and any non-permissive terms
     added in accord with section 7 apply to the code; keep intact all notices
     of the absence of any warranty; and give all recipients a copy of this
     License along with the Program.

     You may charge any price or no price for each copy that you convey, and
     you may offer support or warranty protection for a fee.

     5. Conveying Modified Source Versions.




     You may convey a work based on the Program, or the modifications to
     produce it from the Program, in the form of source code under the terms
     of section 4, provided that you also meet all of these conditions:

     a) The work must carry prominent notices stating that you modified it,
     and giving a relevant date.

     b) The work must carry prominent notices stating that it is released
     under this License and any conditions added under section 7. This
     requirement modifies the requirement in section 4 to "keep intact all
     notices".

     c) You must license the entire work, as a whole, under this License to
     anyone who comes into possession of a copy. This License will therefore
     apply, along with any applicable section 7 additional terms, to the whole
     of the work, and all its parts, regardless of how they are packaged. This
     License gives no permission to license the work in any other way, but it
     does not invalidate such permission if you have separately received it.

     d) If the work has interactive user interfaces, each must display
     Appropriate Legal Notices; however, if the Program has interactive
     interfaces that do not display Appropriate Legal Notices, your work need
     not make them do so.

     A compilation of a covered work with other separate and independent
     works, which are not by their nature extensions of the covered work, and
     which are not combined with it such as to form a larger program, in or on
     a volume of a storage or distribution medium, is called an "aggregate"
     if the compilation and its resulting copyright are not used to limit
     the access or legal rights of the compilation's users beyond what the
     individual works permit. Inclusion of a covered work in an aggregate does
     not cause this License to apply to the other parts of the aggregate.

     6. Conveying Non-Source Forms.

     You may convey a covered work in object code form under the terms of
     sections 4 and 5, provided that you also convey the machine-readable
     Corresponding Source under the terms of this License, in one of these
     ways:

     a) Convey the object code in, or embodied in, a physical product
     (including a physical distribution medium), accompanied by the
     Corresponding Source fixed on a durable physical medium customarily used
     for software interchange.

     b) Convey the object code in, or embodied in, a physical product
     (including a physical distribution medium), accompanied by a written
     offer, valid for at least three years and valid for as long as you




     offer spare parts or customer support for that product model, to
     give anyone who possesses the object code either (1) a copy of the
     Corresponding Source for all the software in the product that is covered
     by this License, on a durable physical medium customarily used for
     software interchange, for a price no more than your reasonable cost of
     physically performing this conveying of source, or (2) access to copy the
     Corresponding Source from a network server at no charge.

     c) Convey individual copies of the object code with a copy of the written
     offer to provide the Corresponding Source. This alternative is allowed
     only occasionally and noncommercially, and only if you received the
     object code with such an offer, in accord with subsection 6b.

     d) Convey the object code by offering access from a designated
     place (gratis or for a charge), and offer equivalent access to the
     Corresponding Source in the same way through the same place at no further
     charge. You need not require recipients to copy the Corresponding Source
     along with the object code. If the place to copy the object code is a
     network server, the Corresponding Source may be on a different server
     (operated by you or a third party) that supports equivalent copying
     facilities, provided you maintain clear directions next to the object
     code saying where to find the Corresponding Source. Regardless of what
     server hosts the Corresponding Source, you remain obligated to ensure
     that it is available for as long as needed to satisfy these requirements.

     e) Convey the object code using peer-to-peer transmission, provided
     you inform other peers where the object code and Corresponding Source
     of the work are being offered to the general public at no charge under
     subsection 6d.

     A separable portion of the object code, whose source code is excluded
     from the Corresponding Source as a System Library, need not be included
     in conveying the object code work.

     A "User Product" is either (1) a "consumer product", which means any
     tangible personal property which is normally used for personal, family,
     or household purposes, or (2) anything designed or sold for incorporation
     into a dwelling. In determining whether a product is a consumer product,
     doubtful cases shall be resolved in favor of coverage. For a particular
     product received by a particular user, "normally used" refers to a
     typical or common use of that class of product, regardless of the
     status of the particular user or of the way in which the particular user
     actually uses, or expects or is expected to use, the product. A product
     is a consumer product regardless of whether the product has substantial
     commercial, industrial or non-consumer uses, unless such uses represent
     the only significant mode of use of the product.




     "Installation Information" for a User Product means any methods,
     procedures, authorization keys, or other information required to install
     and execute modified versions of a covered work in that User Product
     from a modified version of its Corresponding Source. The information
     must suffice to ensure that the continued functioning of the modified
     object code is in no case prevented or interfered with solely because
     modification has been made.

     If you convey an object code work under this section in, or with, or
     specifically for use in, a User Product, and the conveying occurs as
     part of a transaction in which the right of possession and use of the
     User Product is transferred to the recipient in perpetuity or for a
     fixed term (regardless of how the transaction is characterized), the
     Corresponding Source conveyed under this section must be accompanied
     by the Installation Information. But this requirement does not apply if
     neither you nor any third party retains the ability to install modified
     object code on the User Product (for example, the work has been installed
     in ROM).

     The requirement to provide Installation Information does not include a
     requirement to continue to provide support service, warranty, or updates
     for a work that has been modified or installed by the recipient, or
     for the User Product in which it has been modified or installed. Access
     to a network may be denied when the modification itself materially and
     adversely affects the operation of the network or violates the rules and
     protocols for communication across the network.

     Corresponding Source conveyed, and Installation Information provided, in
     accord with this section must be in a format that is publicly documented
     (and with an implementation available to the public in source code form),
     and must require no special password or key for unpacking, reading or
     copying.

     7. Additional Terms.

     "Additional permissions" are terms that supplement the terms of this
     License by making exceptions from one or more of its conditions.
     Additional permissions that are applicable to the entire Program shall
     be treated as though they were included in this License, to the extent
     that they are valid under applicable law. If additional permissions apply
     only to part of the Program, that part may be used separately under those
     permissions, but the entire Program remains governed by this License
     without regard to the additional permissions.

     When you convey a copy of a covered work, you may at your option remove
     any additional permissions from that copy, or from any part of it.
     (Additional permissions may be written to require their own removal




     in certain cases when you modify the work.) You may place additional
     permissions on material, added by you to a covered work, for which you
     have or can give appropriate copyright permission.

     Notwithstanding any other provision of this License, for material you
     add to a covered work, you may (if authorized by the copyright holders of
     that material) supplement the terms of this License with terms:

     a) Disclaiming warranty or limiting liability differently from the terms
     of sections 15 and 16 of this License; or

     b) Requiring preservation of specified reasonable legal notices or
     author attributions in that material or in the Appropriate Legal Notices
     displayed by works containing it; or

     c) Prohibiting misrepresentation of the origin of that material, or
     requiring that modified versions of such material be marked in reasonable
     ways as different from the original version; or

     d) Limiting the use for publicity purposes of names of licensors or
     authors of the material; or

     e) Declining to grant rights under trademark law for use of some trade
     names, trademarks, or service marks; or

     f) Requiring indemnification of licensors and authors of that material
     by anyone who conveys the material (or modified versions of it) with
     contractual assumptions of liability to the recipient, for any liability
     that these contractual assumptions directly impose on those licensors and
     authors.

     All other non-permissive additional terms are considered "further
     restrictions" within the meaning of section 10. If the Program as you
     received it, or any part of it, contains a notice stating that it is
     governed by this License along with a term that is a further restriction,
     you may remove that term. If a license document contains a further
     restriction but permits relicensing or conveying under this License, you
     may add to a covered work material governed by the terms of that license
     document, provided that the further restriction does not survive such
     relicensing or conveying.

     If you add terms to a covered work in accord with this section, you must
     place, in the relevant source files, a statement of the additional terms
     that apply to those files, or a notice indicating where to find the
     applicable terms.




     Additional terms, permissive or non-permissive, may be stated in the
     form of a separately written license, or stated as exceptions; the above
     requirements apply either way.

     8. Termination.

     You may not propagate or modify a covered work except as expressly
     provided under this License. Any attempt otherwise to propagate or modify
     it is void, and will automatically terminate your rights under this
     License (including any patent licenses granted under the third paragraph
     of section 11).

     However, if you cease all violation of this License, then your license
     from a particular copyright holder is reinstated (a) provisionally,
     unless and until the copyright holder explicitly and finally terminates
     your license, and (b) permanently, if the copyright holder fails to
     notify you of the violation by some reasonable means prior to 60 days
     after the cessation.

     Moreover, your license from a particular copyright holder is reinstated
     permanently if the copyright holder notifies you of the violation by
     some reasonable means, this is the first time you have received notice of
     violation of this License (for any work) from that copyright holder, and
     you cure the violation prior to 30 days after your receipt of the notice.

     Termination of your rights under this section does not terminate the
     licenses of parties who have received copies or rights from you under
     this License. If your rights have been terminated and not permanently
     reinstated, you do not qualify to receive new licenses for the same
     material under section 10.

     9. Acceptance Not Required for Having Copies.

     You are not required to accept this License in order to receive or run a
     copy of the Program. Ancillary propagation of a covered work occurring
     solely as a consequence of using peer-to-peer transmission to receive a
     copy likewise does not require acceptance. However, nothing other than
     this License grants you permission to propagate or modify any covered
     work. These actions infringe copyright if you do not accept this License.
     Therefore, by modifying or propagating a covered work, you indicate your
     acceptance of this License to do so.

     10. Automatic Licensing of Downstream Recipients.

     Each time you convey a covered work, the recipient automatically receives
     a license from the original licensors, to run, modify and propagate that
     work, subject to this License. You are not responsible for enforcing
     compliance by third parties with this License.




     An "entity transaction" is a transaction transferring control of an
     organization, or substantially all assets of one, or subdividing an
     organization, or merging organizations. If propagation of a covered work
     results from an entity transaction, each party to that transaction who
     receives a copy of the work also receives whatever licenses to the work
     the party's predecessor in interest had or could give under the previous
     paragraph, plus a right to possession of the Corresponding Source of the
     work from the predecessor in interest, if the predecessor has it or can
     get it with reasonable efforts.

     You may not impose any further restrictions on the exercise of the rights
     granted or affirmed under this License. For example, you may not impose
     a license fee, royalty, or other charge for exercise of rights granted
     under this License, and you may not initiate litigation (including a
     cross-claim or counterclaim in a lawsuit) alleging that any patent claim
     is infringed by making, using, selling, offering for sale, or importing
     the Program or any portion of it.

     11. Patents.

     A "contributor" is a copyright holder who authorizes use under this
     License of the Program or a work on which the Program is based. The work
     thus licensed is called the contributor's "contributor version".

     A contributor's "essential patent claims" are all patent claims owned
     or controlled by the contributor, whether already acquired or hereafter
     acquired, that would be infringed by some manner, permitted by this
     License, of making, using, or selling its contributor version, but do not
     include claims that would be infringed only as a consequence of further
     modification of the contributor version. For purposes of this definition,
     "control" includes the right to grant patent sublicenses in a manner
     consistent with the requirements of this License.

     Each contributor grants you a non-exclusive, worldwide, royalty-free
     patent license under the contributor's essential patent claims, to make,
     use, sell, offer for sale, import and otherwise run, modify and propagate
     the contents of its contributor version.

     In the following three paragraphs, a "patent license" is any express
     agreement or commitment, however denominated, not to enforce a patent
     (such as an express permission to practice a patent or covenant not to
     sue for patent infringement). To "grant" such a patent license to a party
     means to make such an agreement or commitment not to enforce a patent
     against the party.

     If you convey a covered work, knowingly relying on a patent license, and
     the Corresponding Source of the work is not available for anyone to copy,




     free of charge and under the terms of this License, through a publicly
     available network server or other readily accessible means, then you
     must either (1) cause the Corresponding Source to be so available, or
     (2) arrange to deprive yourself of the benefit of the patent license for
     this particular work, or (3) arrange, in a manner consistent with the
     requirements of this License, to extend the patent license to downstream
     recipients. "Knowingly relying" means you have actual knowledge that, but
     for the patent license, your conveying the covered work in a country, or
     your recipient's use of the covered work in a country, would infringe
     one or more identifiable patents in that country that you have reason to
     believe are valid.

     If, pursuant to or in connection with a single transaction or
     arrangement, you convey, or propagate by procuring conveyance of, a
     covered work, and grant a patent license to some of the parties receiving
     the covered work authorizing them to use, propagate, modify or convey a
     specific copy of the covered work, then the patent license you grant is
     automatically extended to all recipients of the covered work and works
     based on it.

     A patent license is "discriminatory" if it does not include within the
     scope of its coverage, prohibits the exercise of, or is conditioned
     on the non-exercise of one or more of the rights that are specifically
     granted under this License. You may not convey a covered work if you
     are a party to an arrangement with a third party that is in the business
     of distributing software, under which you make payment to the third
     party based on the extent of your activity of conveying the work, and
     under which the third party grants, to any of the parties who would
     receive the covered work from you, a discriminatory patent license (a)
     in connection with copies of the covered work conveyed by you (or copies
     made from those copies), or (b) primarily for and in connection with
     specific products or compilations that contain the covered work, unless
     you entered into that arrangement, or that patent license was granted,
     prior to 28 March 2007.

     Nothing in this License shall be construed as excluding or limiting any
     implied license or other defenses to infringement that may otherwise be
     available to you under applicable patent law.

     12. No Surrender of Others' Freedom.

     If conditions are imposed on you (whether by court order, agreement or
     otherwise) that contradict the conditions of this License, they do not
     excuse you from the conditions of this License. If you cannot convey
     a covered work so as to satisfy simultaneously your obligations under
     this License and any other pertinent obligations, then as a consequence




     you may not convey it at all. For example, if you agree to terms that
     obligate you to collect a royalty for further conveying from those to
     whom you convey the Program, the only way you could satisfy both those
     terms and this License would be to refrain entirely from conveying the
     Program.

     13. Use with the GNU Affero General Public License.

     Notwithstanding any other provision of this License, you have permission
     to link or combine any covered work with a work licensed under version
     3 of the GNU Affero General Public License into a single combined
     work, and to convey the resulting work. The terms of this License will
     continue to apply to the part which is the covered work, but the special
     requirements of the GNU Affero General Public License, section 13,
     concerning interaction through a network will apply to the combination as
     such.

     14. Revised Versions of this License.

     The Free Software Foundation may publish revised and/or new versions of
     the GNU General Public License from time to time. Such new versions will
     be similar in spirit to the present version, but may differ in detail to
     address new problems or concerns.

     Each version is given a distinguishing version number. If the Program
     specifies that a certain numbered version of the GNU General Public
     License "or any later version" applies to it, you have the option of
     following the terms and conditions either of that numbered version or
     of any later version published by the Free Software Foundation. If the
     Program does not specify a version number of the GNU General Public
     License, you may choose any version ever published by the Free Software
     Foundation.

     If the Program specifies that a proxy can decide which future versions of
     the GNU General Public License can be used, that proxy's public statement
     of acceptance of a version permanently authorizes you to choose that
     version for the Program.

     Later license versions may give you additional or different permissions.
     However, no additional obligations are imposed on any author or copyright
     holder as a result of your choosing to follow a later version.

     15. Disclaimer of Warranty.

     THERE IS NO WARRANTY FOR THE PROGRAM, TO THE EXTENT PERMITTED BY
     APPLICABLE LAW. EXCEPT WHEN OTHERWISE STATED IN WRITING THE COPYRIGHT
     HOLDERS AND/OR OTHER PARTIES PROVIDE THE PROGRAM "AS IS" WITHOUT WARRANTY
     OF ANY KIND, EITHER EXPRESSED OR IMPLIED, INCLUDING, BUT NOT LIMITED TO,




     THE IMPLIED WARRANTIES OF MERCHANTABILITY AND FITNESS FOR A PARTICULAR
     PURPOSE. THE ENTIRE RISK AS TO THE QUALITY AND PERFORMANCE OF THE PROGRAM
     IS WITH YOU. SHOULD THE PROGRAM PROVE DEFECTIVE, YOU ASSUME THE COST OF
     ALL NECESSARY SERVICING, REPAIR OR CORRECTION.

     16. Limitation of Liability.

     IN NO EVENT UNLESS REQUIRED BY APPLICABLE LAW OR AGREED TO IN WRITING
     WILL ANY COPYRIGHT HOLDER, OR ANY OTHER PARTY WHO MODIFIES AND/OR CONVEYS
     THE PROGRAM AS PERMITTED ABOVE, BE LIABLE TO YOU FOR DAMAGES, INCLUDING
     ANY GENERAL, SPECIAL, INCIDENTAL OR CONSEQUENTIAL DAMAGES ARISING OUT
     OF THE USE OR INABILITY TO USE THE PROGRAM (INCLUDING BUT NOT LIMITED TO
     LOSS OF DATA OR DATA BEING RENDERED INACCURATE OR LOSSES SUSTAINED BY YOU
     OR THIRD PARTIES OR A FAILURE OF THE PROGRAM TO OPERATE WITH ANY OTHER
     PROGRAMS), EVEN IF SUCH HOLDER OR OTHER PARTY HAS BEEN ADVISED OF THE
     POSSIBILITY OF SUCH DAMAGES.

     17. Interpretation of Sections 15 and 16.

     If the disclaimer of warranty and limitation of liability provided above
     cannot be given local legal effect according to their terms, reviewing
     courts shall apply local law that most closely approximates an absolute
     waiver of all civil liability in connection with the Program, unless a
     warranty or assumption of liability accompanies a copy of the Program in
     return for a fee.

     END OF TERMS AND CONDITIONS

     How to Apply These Terms to Your New Programs

     If you develop a new program, and you want it to be of the greatest
     possible use to the public, the best way to achieve this is to make it
     free software which everyone can redistribute and change under these
     terms.

     To do so, attach the following notices to the program. It is safest
     to attach them to the start of each source file to most effectively
     state the exclusion of warranty; and each file should have at least the
     "copyright" line and a pointer to where the full notice is found.

     <one line to give the program's name and a brief idea of what it does.>
     Copyright (C) <year> <name of author>

     This program is free software: you can redistribute it and/or modify it
     under the terms of the GNU General Public License as published by the
     Free Software Foundation, either version 3 of the License, or (at your
     option) any later version.




     This program is distributed in the hope that it will be useful,
     but WITHOUT ANY WARRANTY; without even the implied warranty of
     MERCHANTABILITY or FITNESS FOR A PARTICULAR PURPOSE. See the GNU General
     Public License for more details.

     You should have received a copy of the GNU General Public License along
     with this program. If not, see <http://www.gnu.org/licenses/>.

     Also add information on how to contact you by electronic and paper mail.

     If the program does terminal interaction, make it output a short notice
     like this when it starts in an interactive mode:

     <program> Copyright (C) <year> <name of author> This program comes with
     ABSOLUTELY NO WARRANTY; for details type `show w'. This is free software,
     and you are welcome to redistribute it under certain conditions; type
     `show c' for details.

     The hypothetical commands `show w' and `show c' should show the
     appropriate parts of the General Public License. Of course, your
     program's commands might be different; for a GUI interface, you would use
     an "about box".

     You should also get your employer (if you work as a programmer) or
     school, if any, to sign a "copyright disclaimer" for the program, if
     necessary. For more information on this, and how to apply and follow the
     GNU GPL, see <http://www.gnu.org/licenses/>.

     The GNU General Public License does not permit incorporating your program
     into proprietary programs. If your program is a subroutine library, you
     may consider it more useful to permit linking proprietary applications
     with the library. If this is what you want to do, use the GNU Lesser
     General Public License instead of this License. But first, please read
     <http://www.gnu.org/philosophy/why-not-lgpl.html>.
\end{lstlisting}

